%% ============================================================================
%% Chapter 15: Microservices & Resilient Inter-Service Communication
%% Mastering APIs: From Zero to Production Guru
%% ============================================================================
\chapter{Microservices \& Resilient Inter-Service Communication}
\label{ch:microservices}

%% ----------------------------------------------------------------------------
%% Executive Summary
%% ----------------------------------------------------------------------------
Splitting a monolith into services does not remove coupling; it moves it
from the compiler's reach, where it is cheap and visible, onto the network,
where it is expensive, latent, and partially observable. A method call that
used to take nanoseconds and fail only when the process died now takes
milliseconds, fails in at least five distinguishable ways, and --- this is
the part that hurts --- sometimes \emph{does not fail at all}, it merely
gets slow, and slow is worse than down because slow consumes the caller's
resources while offering no error to react to. This chapter is about
engineering that reality rather than discovering it during an incident.

We begin with the communication taxonomy: synchronous request/response
(REST, gRPC) versus asynchronous messaging (queues, event streams), and
orchestration versus choreography as the two ways to coordinate multi-step
work. We then make coupling precise along three dimensions ---
\emph{temporal} (must both parties be up simultaneously?), \emph{spatial}
(must the caller know where the callee lives?), and \emph{logical} (must
both agree on schema and semantics?) --- and apply the eight fallacies of
distributed computing as a checklist to every API call you write. The
fallacies are not trivia; each one maps to a production failure mode with a
name, a metric, and a mitigation.

The resilience core of the chapter treats five patterns with the rigor they
deserve. Timeouts come first, because every other pattern is built on the
assumption that failure detection is bounded; we derive the deadline budget
rule $T_{\text{caller}} > \sum T_{\text{hops}}$ and show what breaks when it
is violated. The circuit breaker is specified as a formal state machine ---
\textsc{closed}, \textsc{open}, \textsc{half\_open}, with a rolling failure
window, explicit thresholds, and recovery probes --- not as a vibes-based
``it stops calling the broken thing.'' Bulkheads isolate failure domains by
partitioning thread pools, connection pools, and concurrency permits.
Hedged requests attack tail latency and we price their amplification
honestly. Backpressure closes the loop by giving overloaded systems a way
to say no upstream.

Two phenomena get full mathematical treatment. \emph{Retry amplification}:
we prove that an $n$-hop chain in which each hop issues $a$ attempts per
call multiplies one client request into $a\,(a^{n}-1)/(a-1)$ downstream
requests --- with three attempts per hop across three hops, that is $39\times$
load from a single user click (Theorem~\ref{thm:ch15-retry-amplification})
--- and we show how token-bucket retry budgets and jittered, deadline-aware
backoff cap the damage. \emph{Cascading failure}: we walk the death spiral
stage by stage, from a latency blip through pool exhaustion, queue buildup,
and retry storms to fleet-wide outage, with a failure-mode analysis table
that maps each stage to its early signal and its countermeasure.

The architectural second half covers service discovery and load balancing
(client-side, server-side, DNS, registries), the service mesh (sidecar
proxies, mTLS identity, data plane versus control plane --- and a blunt
account of what the mesh does \emph{not} fix, namely your application
semantics), and distributed transactions without two-phase commit: sagas
with compensating actions, the transactional outbox, idempotent consumers,
and how to expose eventual consistency honestly in your API contracts. Every
claim is executable: five typed Python~3.11 listings --- a production-grade
circuit breaker, a retry-budget client, a three-service chaos drill, a
SQLite-backed outbox with relay, and a saga orchestrator --- run fully
offline with deterministic tests, in the Northwind Robotics universe.

\begin{keyconceptbox}
Resilience is a \textbf{budget discipline}, not a library. Every outbound
call spends three budgets: \emph{time} (deadline), \emph{concurrency}
(bulkhead permits), and \emph{retry allowance} (token bucket). A call site
that cannot say how much of each budget it may spend is an unpriced risk.
The circuit breaker then enforces the meta-rule: when a dependency proves
it cannot honor its budget, stop paying until it recovers.
\end{keyconceptbox}

%% ----------------------------------------------------------------------------
\section{Learning Objectives \& Prerequisites}
\label{sec:ch15-objectives}

\subsection*{Learning Objectives}
After completing this chapter, its lab, and its exercises, you will be able
to:

\begin{enumerate}[label=\textbf{LO-\arabic*.}, leftmargin=2.6cm]
  \item Classify any inter-service interaction along the synchronous /
        asynchronous axis and the orchestration / choreography axis, and
        decompose its coupling into temporal, spatial, and logical
        dimensions with a concrete decoupling tactic for each.
  \item Apply the eight fallacies of distributed computing as a design
        review checklist, naming the failure mode and the metric that
        exposes each fallacy when an API call ignores it.
  \item Derive deadline budgets across multi-hop call chains
        ($T_{\text{caller}} > \sum T_{\text{hops}}$), choose timeout values
        from measured latency distributions, and explain why a timeout
        longer than the caller's is worse than no timeout.
  \item Derive the retry-amplification formula
        $S(n, a) = a\,(a^{n}-1)/(a-1)$
        (Theorem~\ref{thm:ch15-retry-amplification}), and design retry
        policies --- budgets, full jitter, deadlines, server-driven
        throttling --- that provably cap amplification.
  \item Specify, implement, and test a circuit breaker as a formal state
        machine: rolling failure window, minimum-call floors, open
        duration, half-open probe limits, and recovery thresholds.
  \item Choose and combine bulkheads, hedged requests, and backpressure
        mechanisms for a given dependency, including the tail-amplification
        cost of hedging.
  \item Trace a cascading failure stage by stage, attach an early-warning
        signal and a countermeasure to each stage, and run a chaos drill
        that measures --- not asserts --- the protection a breaker buys.
  \item Evaluate service discovery and load-balancing topologies
        (client-side, server-side, DNS, registry) and articulate precisely
        which concerns a service mesh absorbs (mTLS, retries, telemetry,
        traffic policy) and which it cannot (idempotency, sagas, schema
        semantics).
  \item Replace distributed transactions with sagas, compensating actions,
        the transactional outbox, and idempotent consumers, and present the
        resulting eventual-consistency boundaries honestly in API contracts.
\end{enumerate}

\subsection*{Prerequisites}
This chapter assumes Chapters~0--14 and leans hardest on four of them:
consuming HTTP APIs in Python, including \texttt{httpx} clients, timeouts,
and transports (Chapter~\ref{ch:consuming_python}); idempotency keys and
concurrency control, which become \emph{load-bearing} once retries and
at-least-once delivery enter the picture (Chapter~\ref{ch:idempotency});
gRPC, whose deadlines, status codes, and streaming flow control recur
throughout (Chapter~\ref{ch:grpc}); and rate limiting, whose token-bucket
machinery we reuse verbatim for retry budgets
(Chapter~\ref{ch:rate_limiting}). Familiarity with building FastAPI
services (Chapter~\ref{ch:fastapi}) and with asynchronous, event-driven API
styles (Chapter~\ref{ch:async_apis}) is assumed in the code sections. All
listings run offline on Python~3.11+; install the dependencies from
Section~\ref{sec:ch15-deps} once, from PowerShell.

%% ----------------------------------------------------------------------------
\section{Deep Technical Mechanics \& Architectural Concepts}
\label{sec:ch15-mechanics}

\subsection{A Taxonomy of Inter-Service Communication}
\label{subsec:ch15-taxonomy}

Every interaction between two services answers four design questions,
whether or not anyone asked them. First, \emph{does the caller block?}
Synchronous communication (REST over HTTP, gRPC) holds the caller's
execution until a response or a timeout arrives; asynchronous communication
(queues, event streams) hands the work to an intermediary and returns.
Second, \emph{who knows the workflow?} In \textbf{orchestration} a
designated coordinator issues commands in sequence and owns the process
state; in \textbf{choreography} each service reacts to events and emits new
ones, and the workflow exists only as the emergent sum of those reactions.
Third, \emph{how many consumers?} Request/response is inherently
point-to-point; event streams are naturally one-to-many. Fourth,
\emph{what is the consistency promise?} Synchronous chains pretend toward
immediacy; asynchronous pipelines admit, and should advertise, eventual
consistency.

\begin{definition}[Orchestration and choreography]
\label{def:ch15-orch-choreo}
A multi-service workflow is \textbf{orchestrated} when one component (the
orchestrator) explicitly invokes each participant in a known order and
tracks overall progress. It is \textbf{choreographed} when no component
knows the whole workflow: each participant subscribes to domain events,
acts, and publishes its own events, and the workflow is the composition of
those local reactions~\cite{newman2021microservices}.
\end{definition}

Neither coordination style is a resilience pattern by itself; they differ in
\emph{where failure is managed}. An orchestrator can retry, compensate, and
report a single workflow status --- at the cost of being a stateful
component you must make highly available. Choreography removes the
coordinator as a bottleneck and a single point of failure, but scatters
workflow state across event logs and inboxes, making ``what is the status
of order ORD-7?'' a query you must deliberately engineer (a read model, a
status API) rather than a row you can \texttt{SELECT}. Production systems
mix both: orchestration inside a bounded context (the order saga of
Listing~\ref{lst:ch15-saga}), choreography across context boundaries
(order-placed events fanning out to analytics, notifications, and fraud
screening).

\begin{table}[h]
  \centering
  \small
  \begin{tabularx}{\textwidth}{l X X}
    \toprule
    \textbf{Dimension} & \textbf{Synchronous (REST/gRPC)} & \textbf{Asynchronous (queue/stream)} \\
    \midrule
    Caller blocked & Yes, until response/timeout & No; hand-off to broker \\
    Temporal coupling & High: both must be up & Low: broker buffers downtime \\
    Failure surface & Timeouts, 5xx, partial response & Duplicates, reordering, lag, poison messages \\
    Latency profile & Additive across hops & Broker hop + consumer pacing \\
    Workflow visibility & Trace IDs across calls & Event correlation, read models \\
    Best for & Queries, user-facing commands & Fan-out, long work, smoothing load \\
    \bottomrule
  \end{tabularx}
  \caption{Synchronous versus asynchronous inter-service communication.}
  \label{tab:ch15-sync-async}
\end{table}

\subsection{Coupling Dimensions and the Fallacies of Distributed Computing}
\label{subsec:ch15-coupling}

``Loosely coupled'' is a slogan until you decompose it. Three dimensions
cover most of what matters between services.

\begin{definition}[Coupling dimensions]
\label{def:ch15-coupling}
Two services are \textbf{temporally coupled} if a call requires both to be
healthy at the same instant; \textbf{spatially coupled} if the caller must
know the callee's network location (host, port, instance identity);
\textbf{logically coupled} if they must agree on schemas, semantics, and
deployment cadence to interoperate.
\end{definition}

Asynchronous messaging attacks temporal coupling directly --- the broker
holds the message while the consumer is down --- and, combined with a
registry, spatial coupling too. Logical coupling is the stubborn one:
schema registries, consumer-driven contracts, and explicit versioning
(Chapter~\ref{ch:versioning}) manage it but never eliminate it, because
semantics live in code on both sides. A useful design-review question for
every integration: \emph{which of the three couplings does this call site
have, and is that deliberate?}

The eight fallacies of distributed computing, first catalogued by Peter
Deutsch and colleagues at Sun Microsystems, read like a list of things
engineers believe precisely until their first major incident. Applied to
inter-service API calls:

\begin{enumerate}
  \item \emph{The network is reliable.} Packets vanish; TCP connections
        half-open; the callee may process your request and lose the
        response. Consequence: you often cannot distinguish ``not done''
        from ``done but unacknowledged'' --- which is why idempotency keys
        (Chapter~\ref{ch:idempotency}) are a prerequisite for safe retries,
        not an optional garnish.
  \item \emph{Latency is zero.} In-process calls cost nanoseconds; network
        calls cost milliseconds, with a tail that is 10--100$\times$ the
        median. Tail latency, not the mean, sets your timeout budget.
  \item \emph{Bandwidth is infinite.} Chatty fine-grained APIs that were
        free inside a process become serialization and egress bills across
        a network. Coarse-grained payloads and pagination
        (Chapter~\ref{ch:pagination}) exist because of this fallacy.
  \item \emph{The network is secure.} Every hop is a trust boundary; mTLS
        and service identity (Section~\ref{subsec:ch15-mesh}) replace the
        fiction of a trusted LAN.
  \item \emph{Topology doesn't change.} Autoscaling, rolling deploys, and
        container rescheduling move your callees constantly; hardcoded
        addresses are latent outages (spatial coupling again).
  \item \emph{There is one administrator.} Your dependency is run by
        another team with its own deploy cadence, SLOs, and incidents;
        version skew is the steady state, not an anomaly.
  \item \emph{Transport cost is zero.} Serialization, encryption, and
        proxying burn CPU on both sides; meshes add hops. Measure the tax.
  \item \emph{The network is homogeneous.} Mixed client libraries, proxy
        generations, and TLS stacks mean the wire contract --- headers,
        status codes, HTTP semantics (Chapter~\ref{ch:http_wire}) --- is
        the only thing you can truly rely on.
\end{enumerate}

\begin{warningbox}
The most dangerous sentence in microservice design is ``it is just an
internal call.'' Internality does not repeal the fallacies; it usually
means \emph{fewer} timeouts, retries, and contracts than the public API
gets. Apply the same discipline inward that you would demand of a
third-party integration --- your VPC is still a network.
\end{warningbox}

\subsection{Timeouts Everywhere, and How to Budget Them}
\label{subsec:ch15-timeouts}

A call without a timeout is a promise to wait forever, and ``forever'' is
the one deadline production will always accept. Timeouts are the foundation
every other pattern in this chapter stands on: the circuit breaker's
failure counter is fed mostly by timeout expirations; bulkheads only help
if blocked workers are eventually released; retry budgets assume each
attempt has a bounded cost.

Timeout selection is a measurement exercise, not a guess. Take the
dependency's latency distribution under realistic load --- not the mean,
the tail --- and set the timeout at a high percentile with headroom, e.g.\
$p99.9$ plus a margin, then validate with the chaos drill of
Section~\ref{sec:ch15-lab}. Too tight and you convert slow successes into
failures (and retries, and amplification); too loose and you hold resources
past the point where the caller still cares.

The structural rule is the \textbf{deadline budget}: a caller's timeout
must exceed the sum of the budgets it grants downstream, or the chain
guarantees wasted work.

\begin{definition}[Deadline budget rule]
\label{def:ch15-deadline}
For a call chain $C_0 \to C_1 \to \dots \to C_n$ where $T_i$ is the timeout
$C_{i-1}$ applies to its call to $C_i$, correct budgeting requires
$T_i > \sum_{j>i} T_j$ for every hop --- equivalently, budgets shrink
strictly toward the leaves. gRPC makes this explicit by propagating a
\emph{deadline} (an absolute timestamp) in metadata, so every hop derives
its remaining budget instead of re-guessing~\cite{grpcdocs}.
\end{definition}

Concretely, for the Northwind Robotics catalog chain: the public API grants
the catalog request a $2\,\mathrm{s}$ end-to-end deadline; catalog budgets
$800\,\mathrm{ms}$ for its inventory call; inventory budgets
$400\,\mathrm{ms}$ for warehouse; warehouse budgets $200\,\mathrm{ms}$ for
its database. Each layer keeps headroom for serialization, queuing, and its
own logic. Section~\ref{sec:ch15-pitfalls} catalogs what happens when a
middle hop violates the rule (a timeout \emph{longer} than the caller's):
the caller abandons the request while the callee keeps working --- the
definition of waste, and under load, of cascading failure.

\subsection{The Mathematics of Retry Amplification}
\label{subsec:ch15-retry-math}

Retries are the first resilience tool everyone reaches for and the first
amplifier of every outage. The amplification is multiplicative across hops,
and it is worth deriving once, carefully, because the numbers are startling
enough to change design instincts.

\begin{theorem}[Retry amplification in an $n$-hop chain]
\label{thm:ch15-retry-amplification}
Consider a call chain of $n$ hops $C_0 \to C_1 \to \dots \to C_n$. Suppose
each service $C_i$ ($i < n$), upon receiving a call, issues up to $a$
attempts to $C_{i+1}$ when attempts fail ($a = 1 + r$ for $r$ retries).
Then one client request at $C_0$ generates up to $a^k$ calls arriving at
hop $k$, and the total number of downstream requests in the system is
\[
  S(n, a) \;=\; \sum_{k=1}^{n} a^{k} \;=\; a\,\frac{a^{n} - 1}{a - 1}.
\]
\end{theorem}

\begin{proof}
By induction on the hop index. One request arrives at $C_0$ (base of the
chain receives the client call). Each call arriving at $C_i$ independently
triggers up to $a$ calls to $C_{i+1}$, so if $L_k$ denotes the number of
calls arriving at hop $k$, then $L_{k+1} = a \cdot L_k$ with $L_0 = 1$,
giving $L_k = a^k$. The system-wide total is the geometric sum
$\sum_{k=1}^{n} a^k = a(a^n - 1)/(a - 1)$.
\end{proof}

The canonical numbers: a three-hop chain (edge $\to$ catalog $\to$
inventory $\to$ warehouse) in which every hop issues up to three attempts
($a = 3$: one initial try plus two retries) produces up to
$3 + 9 + 27 = 39$ downstream requests per client request --- a $39\times$
load multiplier from a single user action, with the leaf service absorbing
$27\times$. If ``three retries'' means three retries \emph{beyond} the
first attempt ($a = 4$), the total is $4 + 16 + 64 = 84$. During an
outage, when failure rates approach $100\%$, these maxima are not worst-case
theory; they are the expected case.

\begin{proposition}[Retry budgets cap amplification]
\label{prop:ch15-budget-cap}
Let each client process hold a token bucket of capacity $B$ where every
retry (not every attempt) consumes one token. Then regardless of configured
\texttt{max\_retries}, the attempts issued per logical request are bounded
by $1 + B$, and the steady-state retry rate is bounded by the refill rate.
Amplification becomes a controlled parameter, not an emergent property.
\end{proposition}

The proof is immediate --- a token must exist for a retry to start --- but
the consequences are deep: the budget converts retries from a per-request
decision into a \emph{per-process allocation problem}, which is the only
granularity at which ``how much extra load may this client add'' is
answerable. Listing~\ref{lst:ch15-retry-budget} implements exactly this,
and Listing~\ref{lst:ch15-retry-budget-tests} proves the cap: configured
for five retries against an always-failing dependency, a client with a
two-token budget issues three requests and stops.

Three more disciplines complete the anti-amplification toolkit.
\textbf{Jitter}: if $N$ clients fail together at time $t$ and all sleep the
same deterministic backoff, they re-arrive together at $t + \Delta$ and
re-collide; \emph{full jitter}, drawing the sleep uniformly from
$[0, \min(c, b \cdot 2^{k})]$, decorrelates the herd while preserving the
exponential envelope. \textbf{Deadline awareness}: a retry whose backoff
would overshoot the caller's remaining deadline is a retry that cannot
help; skip it (Listing~\ref{lst:ch15-retry-budget} raises
\texttt{RetryDeadlineExceededError} rather than donating doomed work
downstream). \textbf{Server-driven throttling}: when the server says 429 or
503 with \texttt{Retry-After}, that signal outranks the client's local
schedule --- the server has information the client lacks, and honoring it
is what turns a thundering herd into an orderly queue
(Chapter~\ref{ch:rate_limiting}).

\subsection{The Circuit Breaker, Specified Formally}
\label{subsec:ch15-breaker}

The circuit breaker, named by Nygard~\cite{nygard2018releaseit}, wraps a
dependency in a state machine that substitutes fast, cheap, honest failure
for slow, expensive, resource-consuming failure. Everyone can draw the
three states; production-grade breakers are distinguished by what the
diagram leaves out: \emph{what counts} as a failure, \emph{how many}
observations are required before the statistics mean anything, \emph{who}
is allowed to test recovery, and \emph{when} old evidence expires.

\begin{definition}[Circuit breaker]
\label{def:ch15-breaker}
A circuit breaker is a per-dependency state machine over
$\{\textsc{closed}, \textsc{open}, \textsc{half\_open}\}$ with parameters
$(\theta, W, N, d, p, s)$: failure-rate threshold $\theta \in (0,1]$,
rolling window $W$ seconds, minimum call count $N$, open duration $d$,
probe concurrency limit $p$, and probe successes $s$ required to close.
Transitions: \textsc{closed} $\to$ \textsc{open} when the window holds at
least $N$ calls with failure rate $\ge \theta$; \textsc{open} $\to$
\textsc{half\_open} when $d$ seconds have elapsed since opening;
\textsc{half\_open} $\to$ \textsc{closed} after $s$ consecutive probe
successes; \textsc{half\_open} $\to$ \textsc{open} on any probe failure.
While \textsc{open}, calls are rejected immediately without touching the
dependency.
\end{definition}

\begin{figure}[h]
\centering
\begin{tikzpicture}[font=\small, >=stealth, node distance=2.9cm]
  \node[flowstep, minimum width=3.1cm, align=center] (closed)
    {\textbf{CLOSED}\\ traffic flows\\ rolling window counts};
  \node[flowstep, minimum width=3.1cm, align=center, right=3.4cm of closed] (open)
    {\textbf{OPEN}\\ reject fast\\ no dependency calls};
  \node[flowstep, minimum width=3.4cm, align=center]
    at ($(closed)!0.5!(open) + (0,-3.0)$) (half)
    {\textbf{HALF\_OPEN}\\ admit $\le p$ probes\\ measure recovery};
  \draw[->, thick] (closed) to[bend left=18]
    node[above, align=center] {failure rate $\ge \theta$\\ over $\ge N$ calls in $W$} (open);
  \draw[->, thick] (open) to[bend left=18]
    node[below, align=center] {$t \ge t_{\text{open}} + d$} (half);
  \draw[->, thick, packtorange] (half) to[bend left=18]
    node[below right=-2pt, align=left] {probe failure\\ (dependency still sick)} (open);
  \draw[->, thick, darkblue] (half) to[bend right=18]
    node[below left=-2pt, align=right] {$s$ consecutive\\ probe successes} (closed);
  \node[note, align=center] at ($(closed) + (0,-1.5)$)
    {window evicts observations older than $W$};
  \node[note, align=center] at ($(open) + (0,-1.5)$)
    {rejection cost $\approx$ 0;\\ caller falls back or 503s};
\end{tikzpicture}
\caption{Circuit breaker state machine with the parameters of
Definition~\ref{def:ch15-breaker}. The half-open state is the recovery
probe mechanism: a strictly limited number of trial calls decide whether
the dependency has healed before normal traffic resumes.}
\label{fig:ch15-breaker-fsm}
\end{figure}

Four specification details do the real work. \textbf{Minimum calls} $N$
prevents statistical noise from tripping the breaker: two failures out of
two calls at 3~a.m.\ must not open a circuit that carries ten thousand
calls a minute at noon. \textbf{The rolling window} $W$ makes the decision
a statement about the \emph{recent} past; a tumbling window would let a
burst of failures at the boundary hide in one window and vanish in the
next. \textbf{Probe limits} $p$ and $s$ make recovery gradual --- one probe
at a time, several successes before closing --- because a dependency that
has just recovered is usually fragile, and a stampede of pent-up traffic is
how recoveries become second outages. \textbf{One breaker per dependency}:
the breaker's state is a hypothesis about \emph{a specific} downstream
resource; sharing one breaker across two dependencies makes a sick
dependency poison a healthy one (Section~\ref{sec:ch15-pitfalls}).

While open, the breaker reduces load on the sick dependency from this
caller to approximately zero --- $p$ probe calls per open duration $d$ ---
which is precisely the headroom the dependency needs to finish recovery.
The caller, meanwhile, must decide what rejection \emph{means}: fail the
request (503 with a problem-details body, Chapter~\ref{ch:problem_details}),
degrade gracefully (serve a cached or partial answer), or queue for later.
The breaker decides \emph{when} to stop calling; the application still owns
\emph{what happens instead}.

\subsection{Bulkheads, Hedged Requests, and Backpressure}
\label{subsec:ch15-bulkheads}

\textbf{Bulkheads} partition resources so that one dependency's failure
cannot consume the capacity needed by everything else --- named, like the
breaker, for shipbuilding, via~\cite{nygard2018releaseit}. In practice a
bulkhead is any of: a dedicated connection pool per dependency (an
\texttt{httpx} client per downstream, with its own \texttt{Limits}), a
semaphore bounding concurrent calls into one dependency (the
\texttt{BULKHEAD\_PERMITS} of Listing~\ref{lst:ch15-demo}), or a separate
worker pool per workload class. The sizing rule: permits for dependency $d$
should be at most what $d$ can healthily serve, and the sum of all permit
pools must leave capacity for everything else the process does --- health
checks, metrics, other endpoints. A bulkhead that is larger than the
process's total capacity is decorative.

\textbf{Hedged requests} attack tail latency: if a call has not answered
within the latency budget's $p95$--$p99$, issue a second, identical call to
a different replica and take whichever answers first. The trade-off is
explicit --- you trade a bounded amount of extra load for a large reduction
in tail latency. If the hedge fires when the first call exceeds delay $h$,
and $\Pr[T > h] = q$, expected load per request becomes $1 + q$; with
$q = 0.05$ that is a 5\% tax for cutting the $p99$ dramatically. The
corollaries are mandatory: hedge only idempotent operations (or carry
idempotency keys), cancel or disregard the loser, and never hedge while the
breaker is open --- hedging into a sick dependency is amplification with a
marketing budget.

\textbf{Backpressure} is the mechanism by which an overloaded component
pushes load back upstream instead of collapsing silently. Its forms:
explicit rejection (429/503 with \texttt{Retry-After}), bounded queues that
fail fast when full (an unbounded queue is not a buffer; it is a latency
generator with a memory leak), credit-based flow control (HTTP/2 stream
windows, gRPC flow control~\cite{grpcdocs}), and consumer-side pacing in
event systems (consumer lag \emph{is} the backpressure signal; Chapter
\ref{ch:async_apis}). A system without backpressure converts overload into
latency for everyone; a system with backpressure converts it into rejection
for some --- which is what makes the difference between a degraded service
and a dead one.

\subsection{Anatomy of a Cascading Failure}
\label{subsec:ch15-cascade}

Cascading failures are the reason this chapter exists. They are also
boringly predictable: the same stages, in the same order, in incident
review after incident review. Figure~\ref{fig:ch15-cascade} shows the
spiral; Table~\ref{tab:ch15-failure-modes} attaches signals and
countermeasures.

\begin{figure}[h]
\centering
\begin{tikzpicture}[font=\small, >=stealth]
  \node[flowstep, minimum width=5.6cm, align=left] (s1)
    {\textbf{1. Latency blip} --- dependency slows (GC, disk, neighbor)};
  \node[flowstep, minimum width=5.6cm, align=left, below=0.35cm of s1] (s2)
    {\textbf{2. Workers block} --- in-flight calls hold threads,\\ connections, memory};
  \node[flowstep, minimum width=5.6cm, align=left, below=0.35cm of s2] (s3)
    {\textbf{3. Pools exhaust} --- no free workers; new requests queue};
  \node[flowstep, minimum width=5.6cm, align=left, below=0.35cm of s3] (s4)
    {\textbf{4. Queues build} --- latency rises further; memory pressure,\\ GC churn; health checks start timing out};
  \node[flowstep, minimum width=5.6cm, align=left, below=0.35cm of s4] (s5)
    {\textbf{5. Timeouts fire, retries multiply} --- $a^{n}$ amplification\\ (Theorem~\ref{thm:ch15-retry-amplification}) hits the sick tier};
  \node[flowstep, minimum width=5.6cm, align=left, below=0.35cm of s5] (s6)
    {\textbf{6. Death spiral} --- callers saturate, load balancers eject\\ and re-add flapping instances; outage spreads upstream};
  \foreach \a/\b in {s1/s2, s2/s3, s3/s4, s4/s5, s5/s6}{
    \draw[->, thick] (\a) -- (\b);
  }
  \draw[->, thick, packtorange, dashed] (s5.west) to[out=180, in=180]
    node[left, align=right, text=packtorange] {feedback:\\ retries make\\ stage 1 worse} (s1.west);
  \node[note, right=0.5cm of s1, align=left] {signal: $p99 \nearrow$, error rate flat};
  \node[note, right=0.5cm of s2, align=left] {signal: in-flight gauge $\nearrow$};
  \node[note, right=0.5cm of s3, align=left] {signal: pool wait time $\nearrow$};
  \node[note, right=0.5cm of s4, align=left] {signal: queue depth, heap, GC time};
  \node[note, right=0.5cm of s5, align=left] {signal: retry rate, downstream $QPS \gg$ client $QPS$};
  \node[note, right=0.5cm of s6, align=left] {signal: multi-service SLO burn};
\end{tikzpicture}
\caption{The cascading-failure death spiral. Each stage feeds the next, and
the retry stage feeds back into the first. Every stage has an observable
early-warning signal and a pattern that breaks the chain
(Table~\ref{tab:ch15-failure-modes}).}
\label{fig:ch15-cascade}
\end{figure}

\begin{table}[h]
  \centering
  \small
  \begin{tabularx}{\textwidth}{l X X l}
    \toprule
    \textbf{Stage} & \textbf{Early signal} & \textbf{Countermeasure} & \textbf{Where} \\
    \midrule
    Latency blip & $p99$ latency rises, errors flat & tight per-hop timeouts & \ref{subsec:ch15-timeouts} \\
    Workers block & in-flight calls gauge climbs & bounded concurrency & \ref{subsec:ch15-bulkheads} \\
    Pools exhaust & pool checkout wait $> 0$ & bulkheads per dependency & \ref{subsec:ch15-bulkheads} \\
    Queues build & queue depth, GC time climb & bounded queues, load shedding & \ref{subsec:ch15-bulkheads} \\
    Retry storm & downstream QPS $\gg$ client QPS & retry budgets, jitter, breakers & \ref{subsec:ch15-retry-math} \\
    Death spiral & correlated multi-service burn & all of the above, drilled & \ref{sec:ch15-lab} \\
    \bottomrule
  \end{tabularx}
  \caption{Failure-mode analysis: each stage of the spiral, its observable
  precursor, and the pattern that breaks the chain.}
  \label{tab:ch15-failure-modes}
\end{table}

Two properties make the spiral treacherous. First, \emph{each stage is
locally rational}: waiting a little longer, queueing a little more, and
retrying once are all defensible decisions in isolation; the catastrophe is
a property of the composition. Second, \emph{recovery is not symmetric with
collapse}: a service that fell over under load $L$ cannot usually be
restarted into load $L$ --- it must be brought back behind a breaker, a
load shedder, or a traffic ramp, or it will fall over again immediately.
The Google SRE book's chapter on cascading failures makes the same point
with considerably larger blast radii; the mechanics are identical at
Northwind Robotics scale.

\subsection{Service Discovery and Load Balancing}
\label{subsec:ch15-discovery}

Resilience patterns assume a \emph{pool} of interchangeable backends behind
each logical dependency; discovery and load balancing are how that pool is
maintained and used. The design space has three axes: who resolves
addresses, who picks an instance, and how membership health is known.

\textbf{Server-side load balancing} puts a dedicated load balancer between
callers and the pool: clients resolve one stable virtual address (usually
via DNS) and the balancer owns health checking, instance selection, and
connection management. Simple for clients, operationally centralized --- and
an extra network hop and a component that must itself be made redundant.
\textbf{Client-side load balancing} gives every client the instance list
(from a registry such as Consul, Eureka, or Kubernetes' API) and lets the
client library pick: fewer hops, per-client policy (weighted, zone-aware,
hedged), at the cost of embedding the logic --- and its bugs --- into every
client, in every language you support. \textbf{DNS-based discovery} is the
lowest common denominator: SRV or A records per service, resolved by the
client's stub resolver. It works everywhere and updates nowhere near fast
enough for rapid membership churn, because TTLs cache staleness into the
system by design.

\begin{table}[h]
  \centering
  \small
  \begin{tabularx}{\textwidth}{l X X X}
    \toprule
    & \textbf{Server-side LB} & \textbf{Client-side LB} & \textbf{DNS-based} \\
    \midrule
    Extra hop & Yes & No & No \\
    Client complexity & None & High & Low \\
    Policy flexibility & Centralized & Per-client, rich & Nearly none \\
    Membership freshness & Fast (active health checks) & Fast (registry watch) & Slow (TTL-bound) \\
    Failure domain & LB tier & Each client & Resolver cache \\
    \bottomrule
  \end{tabularx}
  \caption{Discovery and load-balancing topologies.}
  \label{tab:ch15-discovery}
\end{table}

Whichever topology you choose, the resilience-critical property is
\emph{health-check fidelity}: an instance must be removed from rotation
when it is actually failing to serve (not merely when it fails to respond
to a ping), and re-added gradually, because a freshly started instance hit
with a full share of traffic is stage one of Figure~\ref{fig:ch15-cascade}
all over again.

\subsection{Service Mesh: Data Plane versus Control Plane}
\label{subsec:ch15-mesh}

A service mesh extracts cross-cutting communication concerns --- discovery,
load balancing, mTLS, retries, timeouts, telemetry --- from application code
into infrastructure. The canonical architecture is a fleet of \textbf{sidecar
proxies} (Envoy being the reference implementation) forming the
\textbf{data plane}: every packet in or out of a service transits its local
sidecar, so policy executes one hop from the application, without the
application's cooperation or knowledge. Above them sits the \textbf{control
plane} (Istio, Linkerd's control components), which distributes
configuration --- routes, retry policies, certificates, load-balancing
subsets --- and aggregates telemetry, but never touches a request byte.

\begin{figure}[h]
\centering
\begin{tikzpicture}[font=\small, >=stealth]
  % ---- control plane ----
  \node[flowstep, minimum width=9.6cm, align=center, fill=lightgray] (cp)
    {\textbf{Control plane} (e.g.\ Istio) --- config distribution, certificate authority,\\ policy API, telemetry aggregation. \emph{Not on the request path.}};
  % ---- data plane: two services with sidecars ----
  \node[flowstep, align=center, below=1.6cm of cp, xshift=-2.6cm, minimum width=2.4cm] (appA)
    {\textbf{catalog}\\ (app)};
  \node[flowstep, align=center, right=0.3cm of appA, minimum width=2.4cm, fill=blue!8] (scA)
    {\textbf{sidecar}\\ Envoy};
  \node[flowstep, align=center, right=2.2cm of scA, minimum width=2.4cm, fill=blue!8] (scB)
    {\textbf{sidecar}\\ Envoy};
  \node[flowstep, align=center, right=0.3cm of scB, minimum width=2.4cm] (appB)
    {\textbf{inventory}\\ (app)};
  \node[draw=packtgray, dashed, rounded corners, inner sep=6pt,
        fit=(appA)(scA)(scB)(appB)] (dp) {};
  \node[modcap, below=0.05cm of dp] {data plane --- request path};
  % ---- edges ----
  \draw[<->, thick] (appA) -- (scA);
  \draw[<->, thick] (scB) -- (appB);
  \draw[<->, very thick, darkblue] (scA) --
    node[above, align=center] {mTLS: workload identity\\ retries, timeouts, LB, telemetry}
    node[below, align=center, note] {HTTP/2 (or gRPC) between proxies} (scB);
  \draw[->, dashed, packtgray] (cp.south -| scA) -- (scA.north);
  \draw[->, dashed, packtgray] (cp.south -| scB) -- (scB.north);
  \node[note, rotate=0, anchor=west] at ($(cp.south -| scB) + (0.15,-0.6)$)
    {xDS config, cert rotation};
\end{tikzpicture}
\caption{Service mesh anatomy. The data plane (sidecar proxies) executes
traffic policy on the request path: mTLS identity, retries, timeouts,
load balancing, telemetry. The control plane distributes configuration and
certificates but never touches a request.}
\label{fig:ch15-mesh}
\end{figure}

Three mesh capabilities matter for this chapter. \textbf{mTLS identity}:
each workload receives a short-lived cryptographic identity (a SPIFFE-style
verifiable ID), and the sidecars mutually authenticate and encrypt every
hop --- fallacy four handled by default, with identity-based authorization
(``catalog may call inventory'') replacing network-location-based trust.
\textbf{Data-plane traffic policy}: retries, timeouts, circuit breaking
(Envoy calls it outlier detection plus connection-pool circuit breakers),
and load balancing execute identically for every service regardless of
language, which eliminates the ``every team re-implements resilience
slightly wrong'' problem. \textbf{Uniform telemetry}: golden signals per
hop, with consistent labels, emitted by the proxy rather than by
application diligence.

\begin{warningbox}
The mesh does \textbf{not} fix application semantics. It cannot make your
endpoints idempotent, cannot choose your saga compensations, cannot version
your schemas, cannot know which failures are safe to retry (mesh-level
retries of a non-idempotent POST are exactly as dangerous as client-level
ones), and cannot budget your deadlines end-to-end --- it sees hops, not
user requests. Teams that adopt a mesh and delete Chapter-10--15 discipline
from their services get a beautifully observable, mutually authenticated
cascading failure.
\end{warningbox}

\subsection{Distributed Transactions without 2PC}
\label{subsec:ch15-transactions}

A single database gives you ACID; five services give you a choice of which
letter to sacrifice. Two-phase commit (2PC) --- a coordinator asks all
participants to \emph{prepare}, then commits only if all agree --- is
theoretically tidy and operationally avoided: it is blocking (a coordinator
crash leaves participants holding locks in-doubt until recovery), it
couples the availability of the whole to the availability of every
participant, and its chatty prepare/commit rounds multiply latency across
exactly the network boundaries this chapter has spent its length
mistrusting~\cite{kleppmann2017ddia}. Over HTTP, where participants cannot
even hold a lock-shaped resource across stateless requests, 2PC is an
anti-pattern on arrival (Section~\ref{sec:ch15-pitfalls}).

The pragmatic replacement is the \textbf{saga}: a sequence of local
transactions $T_1, \dots, T_n$, each with a \textbf{compensating action}
$C_i$ that semantically undoes $T_i$ --- \emph{semantically}, because the
world has moved on: you cannot un-send an email, but you can send a
correction; you cannot un-charge a card, but you can refund it.

\begin{definition}[Saga]
\label{def:ch15-saga}
A saga is a sequence of local transactions $T_1, \dots, T_n$ with
compensating actions $C_1, \dots, C_{n-1}$ such that if $T_k$ fails, the
system executes $C_{k-1}, \dots, C_1$ in reverse order. The saga provides
\emph{atomicity of outcome} (all steps complete, or all completed steps are
compensated) without isolation: intermediate states are visible to other
requests between steps.
\end{definition}

Sagas come in the two coordination flavors of
Definition~\ref{def:ch15-orch-choreo}. A \emph{choreographed} saga has each
step emit events that trigger the next participant --- simple for two or
three steps, then increasingly difficult to reason about, because the
compensation graph lives in nobody's code. An \emph{orchestrated} saga puts
the sequence, the state, and the compensation logic in one explicit
component (Listing~\ref{lst:ch15-saga}), which is why orchestration is the
default recommendation once a flow has more than a couple of steps or any
step with a non-trivial compensation.

Two plumbing patterns make sagas reliable rather than aspirational. The
\textbf{transactional outbox} solves the dual-write problem --- ``update my
database \emph{and} emit an event'' cannot be done atomically against two
systems, so you don't: in one local transaction you write the domain row
\emph{and} an outbox row; a separate relay reads the outbox and publishes,
marking each event published only after the broker accepts it
(Listing~\ref{lst:ch15-outbox}). The guarantee is \textbf{at-least-once}
delivery: a crash between publish and mark replays the event, which is why
the second pattern, the \textbf{idempotent consumer}, is not optional ---
the consumer records processed event IDs in the same transaction as its
state change, so replays are absorbed
(Chapter~\ref{ch:idempotency} generalizes this to request scopes).

\begin{figure}[h]
\centering
\begin{tikzpicture}[font=\small, >=stealth]
  % ---- top: orchestrated saga ----
  \node[flowstep, minimum width=2.5cm, align=center] (orch)
    {\textbf{saga}\\ \textbf{orchestrator}};
  \node[flowstep, minimum width=2.6cm, align=center, right=1.5cm of orch] (t1)
    {$T_1$ reserve\\ inventory};
  \node[flowstep, minimum width=2.6cm, align=center, right=1.1cm of t1] (t2)
    {$T_2$ charge\\ payment};
  \node[flowstep, minimum width=2.6cm, align=center, right=1.1cm of t2] (t3)
    {$T_3$ schedule\\ shipment};
  \draw[->, thick] (orch) -- (t1);
  \draw[->, thick] (orch.north east) to[out=30, in=150] (t2.north west);
  \draw[->, thick] (orch.north east) to[out=35, in=150] (t3.north west);
  \draw[->, thick] (t1) -- node[above, note] {ok} (t2);
  \draw[->, thick] (t2) -- node[above, note] {ok} (t3);
  \draw[->, thick, packtorange] (t3.south) to[out=-90, in=-90]
    node[below, align=center, text=packtorange]
    {on failure: compensate in reverse --- $C_2$ refund, $C_1$ release}
    (t1.south);
  \node[modcap, left=0.35cm of orch] {orchestrated saga};
  % ---- bottom: outbox flow ----
  \node[flowstep, align=left, below=2.3cm of orch, minimum width=3.3cm] (svc)
    {\textbf{order service}\\ domain write};
  \node[flowstep, align=left, right=1.1cm of svc, minimum width=3.6cm] (db)
    {\textbf{DB (one tx)}\\ \texttt{orders} + \texttt{outbox} rows};
  \node[flowstep, align=left, right=1.1cm of db, minimum width=2.5cm] (relay)
    {\textbf{relay}\\ poll + publish\\ + mark};
  \node[flowstep, align=left, right=1.1cm of relay, minimum width=2.8cm] (cons)
    {\textbf{consumer}\\ idempotent\\ (dedup table)};
  \draw[->, thick] (svc) -- (db);
  \draw[->, thick] (db) -- node[above, note] {unpublished} (relay);
  \draw[->, thick] (relay) --
    node[above, align=center, note] {at-least-once\\ (crash $\Rightarrow$ replay)} (cons);
  \node[modcap, left=0.35cm of svc] {transactional outbox};
\end{tikzpicture}
\caption{Saga orchestration and the transactional outbox. Top: the
orchestrator drives steps in order and compensates in reverse on failure.
Bottom: the domain write and the outbox insert commit atomically; the relay
publishes at-least-once; the idempotent consumer absorbs replays.}
\label{fig:ch15-saga-outbox}
\end{figure}

\begin{examplebox}
\textbf{Presenting eventual consistency in an API.} A client
\texttt{POST /orders} against the saga system gets
\texttt{202 Accepted} with a \texttt{Location: /orders/ORD-7} header and a
body containing \texttt{"status": "PLACED"} and a
\texttt{"status\_url"} poll target. The order resource exposes its saga
state machine honestly: \texttt{PLACED $\to$ PAID $\to$ SHIPPED}, or
\texttt{COMPENSATED} with a machine-readable \texttt{failure\_reason} and a
\texttt{refund\_status}. What the API must never do is return
\texttt{201 Created} with \texttt{"status": "CONFIRMED"} while shipment is
still pending --- consistency boundaries you hide become incident pages
someone else writes.
\end{examplebox}

%% ----------------------------------------------------------------------------
\section{Production-Ready Code Implementation}
\label{sec:ch15-code}

All five subsystems below run offline on Python~3.11+ against in-process
ASGI transports --- no sockets, no containers, no cloud. Determinism is
structural: time enters through injected clocks, concurrency is bounded by
semaphores and coordinated with events and barriers, and every assertion is
on counts and states, never on wall-clock duration. Install once, from
PowerShell, using the requirements of Section~\ref{sec:ch15-deps}:

\begin{minted}{text}
python -m venv .venv
.\.venv\Scripts\Activate.ps1
pip install -r requirements.txt
pytest -q
python demo_services.py
\end{minted}

\subsection{Circuit Breaker: Full State Machine, Sync and Async}
\label{subsec:ch15-code-breaker}

Listing~\ref{lst:ch15-breaker} implements
Definition~\ref{def:ch15-breaker} without compromise: a rolling
\texttt{deque} of timestamped outcomes with lazy eviction, a minimum-call
floor before the failure rate is even evaluated, lazy
\textsc{open}~$\to$~\textsc{half\_open} transition driven by the injected
clock, bounded concurrent probes, a configurable number of consecutive probe
successes to close, and a metrics hook that emits \texttt{call\_success},
\texttt{call\_failure}, \texttt{rejected}, and \texttt{state\_change}
events --- the hook is how you will watch the drill in
Section~\ref{sec:ch15-lab}. The sync \texttt{call} and async \texttt{acall}
share one state machine; late completions from calls admitted before a
transition are ignored so probe accounting stays exact. Note the deliberate
absence of wall-clock reads: \texttt{SystemClock} is the production default,
and tests inject \texttt{ManualClock}.

\noindent\textbf{Listing 15.1 --- \texttt{circuit\_breaker.py}:
production-grade circuit breaker with injected clock and metrics hooks.}
\label{lst:ch15-breaker}
\begin{minted}{python}
"""Production-grade circuit breaker: full state machine, rolling failure
window, half-open probes, metrics hooks, sync + async support.

Deterministic: time comes from an injected ``Clock``; tests drive a manual
clock. No wall-clock reads inside the state machine.
"""
from __future__ import annotations

import threading
import time
from collections import deque
from collections.abc import Awaitable, Callable
from dataclasses import dataclass
from enum import Enum
from typing import Any, Protocol, TypeVar

T = TypeVar("T")


class Clock(Protocol):
    """Monotonic time source, in seconds."""

    def now(self) -> float: ...


class SystemClock:
    """Production clock backed by ``time.monotonic``."""

    def now(self) -> float:
        return time.monotonic()


class BreakerState(str, Enum):
    CLOSED = "closed"        # traffic flows; failures are counted
    OPEN = "open"            # traffic is rejected without touching the dependency
    HALF_OPEN = "half_open"  # a limited number of probe calls are admitted


class CircuitOpenError(RuntimeError):
    """Raised by ``call``/``acall`` when the breaker refuses the call."""


@dataclass(frozen=True, slots=True)
class BreakerConfig:
    """Tuning knobs for one circuit breaker.

    failure_rate_threshold:      open when failures / calls >= this, in (0, 1]
    window_seconds:              rolling observation window length
    minimum_calls:               window is not evaluated until this many calls
    open_duration:               seconds to stay open before the first probe
    half_open_max_probes:        concurrent probe calls admitted in half-open
    half_open_success_threshold: consecutive probe successes required to close
    """

    failure_rate_threshold: float = 0.5
    window_seconds: float = 30.0
    minimum_calls: int = 10
    open_duration: float = 5.0
    half_open_max_probes: int = 1
    half_open_success_threshold: int = 1

    def __post_init__(self) -> None:
        if not 0.0 < self.failure_rate_threshold <= 1.0:
            raise ValueError("failure_rate_threshold must be in (0, 1]")
        if self.window_seconds <= 0 or self.open_duration <= 0:
            raise ValueError("window_seconds and open_duration must be > 0")
        if self.minimum_calls < 1 or self.half_open_max_probes < 1:
            raise ValueError("minimum_calls and half_open_max_probes must be >= 1")
        if self.half_open_success_threshold < 1:
            raise ValueError("half_open_success_threshold must be >= 1")


MetricsHook = Callable[[str, dict[str, Any]], None]


class CircuitBreaker:
    """A circuit breaker guarding exactly ONE dependency.

    States and transitions (Section 3.5 of the chapter):

        CLOSED  --[window failure rate >= threshold]--> OPEN
        OPEN    --[open_duration elapsed]------------> HALF_OPEN
        HALF_OPEN --[probe success x threshold]------> CLOSED
        HALF_OPEN --[probe failure]------------------> OPEN

    Thread-safe for the sync path; the async path is safe on a single
    event loop. Late completions of calls admitted before a state change
    are ignored, which keeps probe accounting exact.
    """

    def __init__(
        self,
        name: str,
        config: BreakerConfig | None = None,
        *,
        clock: Clock | None = None,
        on_event: MetricsHook | None = None,
        is_failure: Callable[[BaseException], bool] | None = None,
    ) -> None:
        if not name:
            raise ValueError("breaker name is required (one breaker per dependency)")
        self._name = name
        self._config = config or BreakerConfig()
        self._clock: Clock = clock if clock is not None else SystemClock()
        self._on_event = on_event
        # Caller-fault exceptions (e.g. ValueError from validation) can be
        # excluded by supplying is_failure; they count as successes here.
        self._is_failure = is_failure or (lambda exc: True)
        self._state = BreakerState.CLOSED
        self._opened_at = 0.0
        self._window: deque[tuple[float, bool]] = deque()
        self._probes_in_flight = 0
        self._probe_successes = 0
        self._lock = threading.Lock()

    # ---- observability -------------------------------------------------
    @property
    def state(self) -> BreakerState:
        with self._lock:
            self._maybe_half_open(self._clock.now())
            return self._state

    def _emit(self, event: str, **fields: Any) -> None:
        if self._on_event is not None:
            self._on_event(event, {"breaker": self._name, **fields})

    # ---- state machine internals (call with self._lock held) ------------
    def _transition(self, to: BreakerState, now: float) -> None:
        if to is self._state:
            return
        previous = self._state
        self._state = to
        if to is BreakerState.OPEN:
            self._opened_at = now
        if to is BreakerState.HALF_OPEN:
            self._probes_in_flight = 0
            self._probe_successes = 0
        self._emit("state_change", previous=previous.value, current=to.value)

    def _maybe_half_open(self, now: float) -> None:
        if (
            self._state is BreakerState.OPEN
            and now - self._opened_at >= self._config.open_duration
        ):
            self._transition(BreakerState.HALF_OPEN, now)

    def _evict(self, now: float) -> None:
        horizon = now - self._config.window_seconds
        while self._window and self._window[0][0] <= horizon:
            self._window.popleft()

    def _admit(self) -> bool:
        """Return True if the admitted call is a half-open probe."""
        with self._lock:
            now = self._clock.now()
            self._maybe_half_open(now)
            if self._state is BreakerState.OPEN:
                self._emit("rejected", state=self._state.value)
                raise CircuitOpenError(
                    f"circuit {self._name!r} is OPEN; call rejected fast"
                )
            if self._state is BreakerState.HALF_OPEN:
                if self._probes_in_flight >= self._config.half_open_max_probes:
                    self._emit("rejected", state=self._state.value)
                    raise CircuitOpenError(
                        f"circuit {self._name!r} is HALF_OPEN; probe limit reached"
                    )
                self._probes_in_flight += 1
                return True
            return False

    def _after_call(self, *, success: bool, probe: bool) -> None:
        with self._lock:
            now = self._clock.now()
            self._emit(
                "call_success" if success else "call_failure",
                probe=probe,
                state=self._state.value,
            )
            if probe:
                self._probes_in_flight -= 1
                if success:
                    self._probe_successes += 1
                    if (
                        self._probe_successes
                        >= self._config.half_open_success_threshold
                    ):
                        self._window.clear()
                        self._transition(BreakerState.CLOSED, now)
                else:
                    # One failed probe is enough: the dependency is still sick.
                    self._transition(BreakerState.OPEN, now)
                return
            if self._state is not BreakerState.CLOSED:
                return  # admitted before a transition; late completion ignored
            self._window.append((now, success))
            self._evict(now)
            if len(self._window) >= self._config.minimum_calls:
                failures = sum(1 for _, ok in self._window if not ok)
                if failures / len(self._window) >= self._config.failure_rate_threshold:
                    self._transition(BreakerState.OPEN, now)

    # ---- public API ------------------------------------------------------
    def call(self, fn: Callable[..., T], *args: Any, **kwargs: Any) -> T:
        """Run ``fn`` through the breaker (synchronous path)."""
        probe = self._admit()
        try:
            result = fn(*args, **kwargs)
        except BaseException as exc:
            self._after_call(success=not self._is_failure(exc), probe=probe)
            raise
        self._after_call(success=True, probe=probe)
        return result

    async def acall(self, fn: Callable[..., Awaitable[T]], *args: Any, **kwargs: Any) -> T:
        """Run an awaitable through the breaker (asyncio path)."""
        probe = self._admit()
        try:
            result = await fn(*args, **kwargs)
        except BaseException as exc:
            self._after_call(success=not self._is_failure(exc), probe=probe)
            raise
        self._after_call(success=True, probe=probe)
        return result
\end{minted}

Listing~\ref{lst:ch15-breaker-tests} proves the documented behavior with a
manual clock: threshold opening at exactly $N$ calls, fail-fast rejection
(the dependency is never invoked while open --- asserted via a counting
stub), the probe protocol (early probe rejected, success closes, failure
reopens and restarts the open duration), the single-concurrent-probe limit
enforced with a real thread blocked on an event, rolling-window eviction,
caller-fault classification, event ordering, and the async path.

\noindent\textbf{Listing 15.2 --- \texttt{test\_circuit\_breaker.py}:
deterministic state-machine proofs.}
\label{lst:ch15-breaker-tests}
\begin{minted}{python}
"""Deterministic tests for the circuit breaker (injected manual clock)."""
from __future__ import annotations

import asyncio
import threading

import pytest

from circuit_breaker import (
    BreakerConfig,
    BreakerState,
    CircuitBreaker,
    CircuitOpenError,
)


class ManualClock:
    """Virtual clock: nothing happens unless the test advances it."""

    def __init__(self, start: float = 1_000.0) -> None:
        self._now = start

    def now(self) -> float:
        return self._now

    def advance(self, seconds: float) -> None:
        if seconds < 0:
            raise ValueError("time only moves forward")
        self._now += seconds


def make_breaker(**overrides) -> tuple[CircuitBreaker, ManualClock, list]:
    clock = ManualClock()
    events: list[tuple[str, dict]] = []
    config = BreakerConfig(
        failure_rate_threshold=overrides.get("failure_rate_threshold", 0.5),
        window_seconds=overrides.get("window_seconds", 10.0),
        minimum_calls=overrides.get("minimum_calls", 4),
        open_duration=overrides.get("open_duration", 5.0),
        half_open_max_probes=overrides.get("half_open_max_probes", 1),
        half_open_success_threshold=overrides.get("half_open_success_threshold", 1),
    )
    breaker = CircuitBreaker(
        "inventory",
        config,
        clock=clock,
        on_event=lambda event, fields: events.append((event, fields)),
    )
    return breaker, clock, events


def boom() -> None:
    raise ConnectionError("inventory unreachable")


def test_closed_breaker_passes_calls_through() -> None:
    breaker, _, _ = make_breaker()
    assert breaker.state is BreakerState.CLOSED
    assert breaker.call(lambda: 42) == 42


def test_opens_when_failure_rate_threshold_is_reached() -> None:
    breaker, _, events = make_breaker()
    # 4 calls, 2 failures -> exactly 0.5 failure rate at minimum_calls.
    breaker.call(lambda: "ok")
    with pytest.raises(ConnectionError):
        breaker.call(boom)
    with pytest.raises(ConnectionError):
        breaker.call(boom)
    assert breaker.state is BreakerState.CLOSED  # 2/3 < window minimum
    breaker.call(lambda: "ok")
    assert breaker.state is BreakerState.OPEN  # 2/4 = 0.5 -> open
    assert any(e == "state_change" and f["current"] == "open" for e, f in events)


def test_open_breaker_rejects_without_calling_the_dependency() -> None:
    breaker, _, _ = make_breaker()
    for _ in range(2):
        breaker.call(lambda: "ok")
        with pytest.raises(ConnectionError):
            breaker.call(boom)
    assert breaker.state is BreakerState.OPEN
    calls = 0

    def counted() -> int:
        nonlocal calls
        calls += 1
        return calls

    for _ in range(10):
        with pytest.raises(CircuitOpenError):
            breaker.call(counted)
    assert calls == 0  # fail-fast: the dependency was never touched


def test_half_open_probe_success_closes_breaker() -> None:
    breaker, clock, _ = make_breaker()
    for _ in range(2):
        with pytest.raises(ConnectionError):
            breaker.call(boom)
        with pytest.raises(ConnectionError):
            breaker.call(boom)
    assert breaker.state is BreakerState.OPEN
    clock.advance(4.9)
    with pytest.raises(CircuitOpenError):
        breaker.call(lambda: "too early")
    clock.advance(0.2)  # now past open_duration
    assert breaker.state is BreakerState.HALF_OPEN
    assert breaker.call(lambda: "recovered") == "recovered"
    assert breaker.state is BreakerState.CLOSED


def test_half_open_probe_failure_reopens() -> None:
    breaker, clock, _ = make_breaker(minimum_calls=2, failure_rate_threshold=1.0)
    for _ in range(2):
        with pytest.raises(ConnectionError):
            breaker.call(boom)
    assert breaker.state is BreakerState.OPEN
    clock.advance(5.0)
    with pytest.raises(ConnectionError):
        breaker.call(boom)  # probe fails
    assert breaker.state is BreakerState.OPEN
    # Reopened: the full open_duration restarts from now.
    clock.advance(4.9)
    with pytest.raises(CircuitOpenError):
        breaker.call(lambda: "still closed to traffic")


def test_half_open_admits_only_one_concurrent_probe() -> None:
    breaker, clock, _ = make_breaker(minimum_calls=2, failure_rate_threshold=1.0)
    for _ in range(2):
        with pytest.raises(ConnectionError):
            breaker.call(boom)
    clock.advance(5.0)

    release = threading.Event()
    probe_started = threading.Event()

    def slow_probe() -> str:
        probe_started.set()
        assert release.wait(timeout=5.0)
        return "healed"

    result: list[str] = []
    worker = threading.Thread(
        target=lambda: result.append(breaker.call(slow_probe)), daemon=True
    )
    worker.start()
    assert probe_started.wait(timeout=5.0)  # probe is in flight
    with pytest.raises(CircuitOpenError):
        breaker.call(lambda: "second probe")  # probe limit = 1
    release.set()
    worker.join(timeout=5.0)
    assert result == ["healed"]
    assert breaker.state is BreakerState.CLOSED


def test_rolling_window_ages_out_old_failures() -> None:
    breaker, clock, _ = make_breaker(window_seconds=10.0, minimum_calls=4)
    # Two failures, then the window slides past them.
    for _ in range(2):
        with pytest.raises(ConnectionError):
            breaker.call(boom)
    clock.advance(11.0)
    breaker.call(lambda: "ok")
    breaker.call(lambda: "ok")
    breaker.call(lambda: "ok")
    breaker.call(lambda: "ok")
    assert breaker.state is BreakerState.CLOSED  # failures evicted: 0/4 rate


def test_caller_fault_exceptions_count_as_successes() -> None:
    clock = ManualClock()
    breaker = CircuitBreaker(
        "catalog",
        BreakerConfig(minimum_calls=2, failure_rate_threshold=0.5),
        clock=clock,
        is_failure=lambda exc: isinstance(exc, ConnectionError),
    )
    for _ in range(5):
        with pytest.raises(ValueError):
            breaker.call(lambda: (_ for _ in ()).throw(ValueError("bad input")))
    assert breaker.state is BreakerState.CLOSED


def test_metrics_hook_receives_events_in_order() -> None:
    breaker, _, events = make_breaker(minimum_calls=2, failure_rate_threshold=1.0)
    with pytest.raises(ConnectionError):
        breaker.call(boom)
    with pytest.raises(ConnectionError):
        breaker.call(boom)
    with pytest.raises(CircuitOpenError):
        breaker.call(lambda: "rejected")
    names = [name for name, _ in events]
    assert names == ["call_failure", "call_failure", "state_change", "rejected"]


def test_async_path_uses_the_same_state_machine() -> None:
    clock = ManualClock()

    async def main() -> None:
        breaker = CircuitBreaker(
            "warehouse",
            BreakerConfig(minimum_calls=2, failure_rate_threshold=1.0),
            clock=clock,
        )

        async def ok() -> str:
            return "fine"

        async def bad() -> str:
            raise ConnectionError("down")

        with pytest.raises(ConnectionError):
            await breaker.acall(bad)
        with pytest.raises(ConnectionError):
            await breaker.acall(bad)
        assert breaker.state is BreakerState.OPEN
        with pytest.raises(CircuitOpenError):
            await breaker.acall(ok)
        clock.advance(5.0)
        assert await breaker.acall(ok) == "fine"
        assert breaker.state is BreakerState.CLOSED

    asyncio.run(main())
\end{minted}

\subsection{Retry-Budget Client: Token Bucket, Full Jitter, Deadlines}
\label{subsec:ch15-code-retry}

Listing~\ref{lst:ch15-retry-budget} turns
Proposition~\ref{prop:ch15-budget-cap} into code. The first attempt is
always free; each retry spends one token from a process-wide bucket; the
bucket refills at a configured rate, so \emph{sustained} retry traffic is
rate-limited even when \emph{bursts} are tolerated. Backoff ceilings grow
exponentially and are jittered by an injectable jitter function (the default
production choice is full jitter over a seeded \texttt{random.Random});
\texttt{no\_jitter} exists for deterministic tests. The deadline rule of
Definition~\ref{def:ch15-deadline} is enforced client-side: if the next
backoff would overshoot the call deadline, the client raises
\texttt{RetryDeadlineExceededError} instead of donating doomed work
downstream.

\noindent\textbf{Listing 15.3 --- \texttt{retry\_budget.py}:
budgeted, jittered, deadline-aware retry client.}
\label{lst:ch15-retry-budget}
\begin{minted}{python}
"""Retry-budgeted HTTP client: token-bucket retry budget + jittered
backoff + deadline awareness.

Design rules implemented (Section 3.4):
  * The first attempt is free; every RETRY spends one budget token.
  * Tokens refill at a fixed rate, so steady-state retry traffic cannot
    exceed ``budget_refill_per_second`` retries/s process-wide.
  * Backoff is exponential with full jitter from a seeded RNG.
  * A call-level deadline caps total elapsed (virtual) time; a retry that
    would overshoot the deadline is never attempted.
"""
from __future__ import annotations

import random
import time
from collections.abc import Callable
from dataclasses import dataclass, field
from typing import Any, Protocol

import httpx


class Clock(Protocol):
    def now(self) -> float: ...
    def sleep(self, seconds: float) -> None: ...


class SystemClock:
    def now(self) -> float:
        return time.monotonic()

    def sleep(self, seconds: float) -> None:
        time.sleep(seconds)


class ManualClock:
    """Virtual clock for deterministic tests; ``sleep`` advances it."""

    def __init__(self, start: float = 0.0) -> None:
        self._now = start

    def now(self) -> float:
        return self._now

    def sleep(self, seconds: float) -> None:
        if seconds < 0:
            raise ValueError("cannot sleep for a negative duration")
        self._now += seconds


class RetryBudgetExhaustedError(RuntimeError):
    """The token bucket is empty: this retry would exceed the retry budget."""


class RetryDeadlineExceededError(RuntimeError):
    """The next backoff sleep would overshoot the call deadline."""


@dataclass(frozen=True, slots=True)
class RetryPolicy:
    max_retries: int = 3                      # retries, not attempts
    base_backoff: float = 0.05                # seconds; delay = base * 2**attempt
    max_backoff: float = 1.0
    deadline_seconds: float = 5.0             # whole-call deadline
    budget_tokens: float = 10.0               # bucket capacity (burst)
    budget_refill_per_second: float = 2.0     # sustained retry allowance
    retry_statuses: frozenset[int] = frozenset({429, 502, 503, 504})

    def __post_init__(self) -> None:
        if self.max_retries < 0:
            raise ValueError("max_retries must be >= 0")
        if self.base_backoff <= 0 or self.max_backoff < self.base_backoff:
            raise ValueError("require 0 < base_backoff <= max_backoff")
        if self.deadline_seconds <= 0:
            raise ValueError("deadline_seconds must be > 0")
        if self.budget_tokens < 0 or self.budget_refill_per_second <= 0:
            raise ValueError("budget parameters must be non-negative/positive")


class TokenBucketRetryBudget:
    """Process-wide retry budget. Tokens are spent ONLY by retries."""

    def __init__(self, capacity: float, refill_per_second: float, clock: Clock) -> None:
        if capacity <= 0 or refill_per_second <= 0:
            raise ValueError("capacity and refill_per_second must be > 0")
        self._capacity = capacity
        self._rate = refill_per_second
        self._clock = clock
        self._tokens = capacity
        self._last = clock.now()

    def _refill(self) -> None:
        now = self._clock.now()
        self._tokens = min(self._capacity, self._tokens + (now - self._last) * self._rate)
        self._last = now

    def consume(self) -> bool:
        self._refill()
        if self._tokens >= 1.0:
            self._tokens -= 1.0
            return True
        return False

    @property
    def tokens(self) -> float:
        self._refill()
        return self._tokens


# Full jitter (AWS Architecture Blog style): uniform in [0, computed cap].
def full_jitter(rng: random.Random) -> Callable[[float], float]:
    def apply(ceiling: float) -> float:
        return rng.uniform(0.0, ceiling)

    return apply


def no_jitter(ceiling: float) -> float:
    return ceiling


@dataclass(slots=True)
class RetryBudgetClient:
    """Wraps an ``httpx.Client``-shaped object with budgeted retries."""

    client: Any  # httpx.Client, or a stub with .request(method, url, **kw)
    budget: TokenBucketRetryBudget
    clock: Clock
    policy: RetryPolicy = field(default_factory=RetryPolicy)
    jitter: Callable[[float], float] = no_jitter
    on_event: Callable[[str, dict[str, Any]], None] | None = None

    def _emit(self, event: str, **fields: Any) -> None:
        if self.on_event is not None:
            self.on_event(event, fields)

    def backoff_ceiling(self, attempt: int) -> float:
        return min(self.policy.max_backoff, self.policy.base_backoff * (2**attempt))

    def request(self, method: str, url: str, **kwargs: Any) -> httpx.Response:
        deadline = self.clock.now() + self.policy.deadline_seconds
        attempt = 0
        while True:
            self._emit("attempt", attempt=attempt, url=url)
            response: httpx.Response | None
            last_error: httpx.TransportError | None = None
            try:
                response = self.client.request(method, url, **kwargs)
            except httpx.TransportError as exc:
                response = None
                last_error = exc
                retryable = True
            else:
                retryable = response.status_code in self.policy.retry_statuses
                if not retryable:
                    return response

            if attempt >= self.policy.max_retries:
                self._emit("retries_exhausted", attempts=attempt + 1, url=url)
                if response is not None:
                    return response  # caller sees the final 5xx
                raise last_error  # type: ignore[misc]

            if not self.budget.consume():
                self._emit("budget_exhausted", attempts=attempt + 1, url=url)
                raise RetryBudgetExhaustedError(
                    "retry budget exhausted; failing fast instead of amplifying"
                ) from last_error

            delay = self.jitter(self.backoff_ceiling(attempt))
            if self.clock.now() + delay > deadline:
                self._emit("deadline_exceeded", attempts=attempt + 1, url=url)
                raise RetryDeadlineExceededError(
                    f"next backoff of {delay:.3f}s would overshoot the deadline"
                ) from last_error

            self._emit("retry", attempt=attempt + 1, delay=delay, url=url)
            self.clock.sleep(delay)
            attempt += 1
\end{minted}

The tests in Listing~\ref{lst:ch15-retry-budget-tests} are the executable
core of this chapter's anti-retry-storm argument:
\texttt{test\_retry\_budget\_caps\_amplification} configures five permitted
retries against an always-failing dependency with a two-token budget and
asserts the client issues exactly three requests --- the budget, not the
retry count, is the binding constraint.

\noindent\textbf{Listing 15.4 --- \texttt{test\_retry\_budget.py}:
budget-cap, refill, deadline, and jitter determinism proofs.}
\label{lst:ch15-retry-budget-tests}
\begin{minted}{python}
"""Deterministic tests for the retry-budget client.

The amplification test is the heart of the chapter's anti-retry-storm
argument: with a budget of 2 tokens, a client configured for up to 5
retries issues AT MOST 3 requests against an always-failing dependency.
"""
from __future__ import annotations

import random

import httpx
import pytest

from retry_budget import (
    ManualClock,
    RetryBudgetClient,
    RetryBudgetExhaustedError,
    RetryDeadlineExceededError,
    RetryPolicy,
    TokenBucketRetryBudget,
    full_jitter,
    no_jitter,
)


class StubClient:
    """httpx.Client-shaped stub; serves queued statuses or transport errors."""

    def __init__(self, script: list[int | type[Exception]]) -> None:
        self._script = list(script)
        self.calls = 0

    def request(self, method: str, url: str, **kwargs) -> httpx.Response:
        self.calls += 1
        outcome = self._script.pop(0) if self._script else 200
        if isinstance(outcome, type) and issubclass(outcome, Exception):
            raise outcome("injected transport failure", request=None)
        return httpx.Response(outcome, request=httpx.Request(method, url))


def make_client(
    script: list[int | type[Exception]],
    *,
    tokens: float = 10.0,
    policy: RetryPolicy | None = None,
    jitter=no_jitter,
    clock: ManualClock | None = None,
) -> tuple[RetryBudgetClient, StubClient, ManualClock, list]:
    clock = clock or ManualClock()
    stub = StubClient(script)
    events: list[tuple[str, dict]] = []
    client = RetryBudgetClient(
        client=stub,
        budget=TokenBucketRetryBudget(tokens, refill_per_second=1.0, clock=clock),
        clock=clock,
        policy=policy or RetryPolicy(max_retries=5, deadline_seconds=60.0),
        jitter=jitter,
        on_event=lambda event, fields: events.append((event, fields)),
    )
    return client, stub, clock, events


def test_success_first_try_spends_no_budget() -> None:
    client, stub, _, _ = make_client([200], tokens=2.0)
    response = client.request("GET", "http://inventory.test/availability/SKU-1")
    assert response.status_code == 200
    assert stub.calls == 1
    assert client.budget.tokens == 2.0


def test_transient_failures_then_success() -> None:
    client, stub, _, events = make_client([503, 503, 200], tokens=5.0)
    response = client.request("GET", "http://inventory.test/availability/SKU-1")
    assert response.status_code == 200
    assert stub.calls == 3
    retries = [f for e, f in events if e == "retry"]
    assert [r["attempt"] for r in retries] == [1, 2]
    assert retries[0]["delay"] == pytest.approx(0.05)
    assert retries[1]["delay"] == pytest.approx(0.10)


def test_retry_budget_caps_amplification() -> None:
    """5 allowed retries, budget of 2 tokens -> at most 3 requests, then fail fast."""
    client, stub, _, events = make_client(
        [503] * 100,
        tokens=2.0,
        policy=RetryPolicy(max_retries=5, deadline_seconds=60.0),
    )
    with pytest.raises(RetryBudgetExhaustedError):
        client.request("GET", "http://inventory.test/availability/SKU-1")
    assert stub.calls == 3  # 1 original + 2 budgeted retries, NOT 6
    assert any(e == "budget_exhausted" for e, _ in events)


def test_budget_refills_over_virtual_time() -> None:
    clock = ManualClock()
    client, stub, _, _ = make_client(
        [503] * 100, tokens=1.0, clock=clock,
        policy=RetryPolicy(max_retries=5, deadline_seconds=60.0),
    )
    with pytest.raises(RetryBudgetExhaustedError):
        client.request("GET", "http://inventory.test/x")
    assert stub.calls == 2
    clock.sleep(2.0)  # refill 1 token/s -> 2 attempts possible again
    with pytest.raises(RetryBudgetExhaustedError):
        client.request("GET", "http://inventory.test/x")
    assert stub.calls == 4


def test_deadline_stops_retries_before_overshoot() -> None:
    # base 0.05s backoff; deadline 0.07s: first retry (0.05) fits,
    # second retry's 0.10s sleep would overshoot -> deadline error.
    client, stub, _, events = make_client(
        [503] * 100,
        tokens=10.0,
        policy=RetryPolicy(max_retries=5, base_backoff=0.05, deadline_seconds=0.07),
    )
    with pytest.raises(RetryDeadlineExceededError):
        client.request("GET", "http://inventory.test/availability/SKU-1")
    assert stub.calls == 2
    assert any(e == "deadline_exceeded" for e, _ in events)


def test_non_retryable_status_returns_immediately() -> None:
    client, stub, _, _ = make_client([404], tokens=5.0)
    response = client.request("GET", "http://inventory.test/missing")
    assert response.status_code == 404
    assert stub.calls == 1
    assert client.budget.tokens == 5.0


def test_transport_errors_are_retryable_and_final_error_reraises() -> None:
    client, stub, _, _ = make_client(
        [httpx.ConnectError, httpx.ConnectError],
        tokens=5.0,
        policy=RetryPolicy(max_retries=1, deadline_seconds=60.0),
    )
    with pytest.raises(httpx.ConnectError):
        client.request("GET", "http://inventory.test/x")
    assert stub.calls == 2  # 1 original + the single permitted retry


def test_full_jitter_is_deterministic_under_a_seeded_rng() -> None:
    rng_a = random.Random(42)
    rng_b = random.Random(42)
    jitter_a = full_jitter(rng_a)
    jitter_b = full_jitter(rng_b)
    ceilings = [0.05, 0.10, 0.20]
    draws_a = [jitter_a(c) for c in ceilings]
    draws_b = [jitter_b(c) for c in ceilings]
    assert draws_a == draws_b
    for ceiling, draw in zip(ceilings, draws_a):
        assert 0.0 <= draw <= ceiling
\end{minted}

\subsection{Three-Service Chaos Drill: Catalog $\to$ Inventory $\to$ Warehouse}
\label{subsec:ch15-code-demo}

Listing~\ref{lst:ch15-demo} wires three FastAPI services --- catalog,
inventory, warehouse --- with \texttt{httpx.AsyncClient} over
\texttt{ASGITransport}: real ASGI request/response cycles and real
\texttt{httpx} timeout machinery, zero network. The catalog edge owns the
resilience posture: a per-hop timeout tighter than its caller's budget, a
dedicated client with a bounded connection pool (bulkhead one), a semaphore
bounding concurrent calls into inventory (bulkhead two), and --- switched
by a flag so the drill can measure the difference --- the
Listing~\ref{lst:ch15-breaker} circuit breaker. The \texttt{ChaosConfig}
knob flips inventory into a hard-down failure mode; \texttt{run\_drill}
then fires thirty concurrent catalog requests and counts what happens.

\noindent\textbf{Listing 15.5 --- \texttt{demo\_services.py}:
three in-process services plus a chaos harness with before/after
measurement.}
\label{lst:ch15-demo}
\begin{minted}{python}
"""Three-service Northwind Robotics demo: catalog -> inventory -> warehouse.

All services are in-process FastAPI apps wired together with
``httpx.AsyncClient`` over ``ASGITransport``: real ASGI request/response
cycles, real timeouts, zero sockets, zero network.

Resilience posture at the catalog edge:
  * per-hop timeouts (catalog->inventory tighter than the caller's budget),
  * a circuit breaker around the inventory dependency (optional, so the
    chaos drill can measure before/after),
  * a bulkhead: a semaphore bounding concurrent calls into inventory.

The chaos harness flips the inventory hop into a failure or latency mode
and measures how many requests actually reach it, with and without the
breaker.
"""
from __future__ import annotations

import asyncio
import time
from dataclasses import dataclass, field
from typing import Any

import httpx
from fastapi import FastAPI, Request
from fastapi.responses import JSONResponse

from circuit_breaker import (
    BreakerConfig,
    CircuitBreaker,
    CircuitOpenError,
)

CATALOG_TIMEOUT = 0.25      # seconds; the caller's budget is larger
INVENTORY_TIMEOUT = 0.20
WAREHOUSE_TIMEOUT = 0.15
BULKHEAD_PERMITS = 8        # max concurrent catalog -> inventory calls

CATALOG = {
    "NR-1001": {"sku": "NR-1001", "name": "Warehouse Picker Bot", "price": 18499.00},
    "NR-2002": {"sku": "NR-2002", "name": "Conveyor Vision Kit", "price": 3299.50},
}


@dataclass
class ChaosConfig:
    """Mutable fault-injection knobs for the inventory service."""

    always_fail: bool = False
    latency_seconds: float = 0.0


@dataclass
class ServiceStats:
    inventory_calls: int = 0
    warehouse_calls: int = 0


# --------------------------------------------------------------------------
# Warehouse (leaf service)
# --------------------------------------------------------------------------
def build_warehouse_app(stats: ServiceStats) -> FastAPI:
    app = FastAPI(title="warehouse")

    @app.get("/stock/{sku}")
    async def stock(sku: str) -> dict[str, Any]:
        stats.warehouse_calls += 1
        return {"sku": sku, "on_hand": 7, "reserved": 2}

    return app


# --------------------------------------------------------------------------
# Inventory (middle service; chaos target)
# --------------------------------------------------------------------------
def build_inventory_app(
    warehouse_client: httpx.AsyncClient,
    chaos: ChaosConfig,
    stats: ServiceStats,
) -> FastAPI:
    app = FastAPI(title="inventory")

    @app.get("/availability/{sku}")
    async def availability(sku: str) -> dict[str, Any]:
        stats.inventory_calls += 1
        if chaos.latency_seconds > 0:
            await asyncio.sleep(chaos.latency_seconds)
        if chaos.always_fail:
            return JSONResponse(
                status_code=500,
                content={"detail": "chaos: inventory database unreachable"},
            )
        stock = await warehouse_client.get(f"/stock/{sku}", timeout=WAREHOUSE_TIMEOUT)
        body = stock.json()
        return {
            "sku": sku,
            "available": body["on_hand"] - body["reserved"],
        }

    return app


# --------------------------------------------------------------------------
# Catalog (edge service; owns the resilience policies)
# --------------------------------------------------------------------------
def build_catalog_app(
    inventory_client: httpx.AsyncClient,
    breaker: CircuitBreaker | None,
    bulkhead: asyncio.Semaphore,
) -> FastAPI:
    app = FastAPI(title="catalog")

    @app.get("/products/{sku}")
    async def product(sku: str, request: Request) -> Any:
        product_row = CATALOG.get(sku)
        if product_row is None:
            return JSONResponse(status_code=404, content={"detail": "unknown sku"})

        async def fetch_availability() -> dict[str, Any]:
            async with bulkhead:  # bulkhead: bounded concurrency into inventory
                response = await inventory_client.get(
                    f"/availability/{sku}", timeout=CATALOG_TIMEOUT
                )
            response.raise_for_status()
            return response.json()

        try:
            if breaker is not None:
                availability = await breaker.acall(fetch_availability)
            else:
                availability = await fetch_availability()
        except CircuitOpenError:
            return JSONResponse(
                status_code=503,
                content={"detail": "inventory unavailable (circuit open)",
                         "circuit_open": True},
            )
        except (httpx.HTTPError, asyncio.TimeoutError) as exc:
            return JSONResponse(
                status_code=503,
                content={"detail": f"inventory call failed: {type(exc).__name__}",
                         "circuit_open": False},
            )
        return {**product_row, "availability": availability}

    return app


# --------------------------------------------------------------------------
# Wiring + chaos drill harness
# --------------------------------------------------------------------------
@dataclass
class DrillResult:
    requests: int
    catalog_errors: int
    circuit_open_rejections: int
    inventory_calls: int
    warehouse_calls: int
    wall_seconds: float


@dataclass
class DemoSystem:
    apps: dict[str, FastAPI]
    chaos: ChaosConfig
    stats: ServiceStats
    events: list[tuple[str, dict]] = field(default_factory=list)


def build_system(*, use_breaker: bool) -> DemoSystem:
    chaos = ChaosConfig()
    stats = ServiceStats()
    events: list[tuple[str, dict]] = []

    warehouse_app = build_warehouse_app(stats)
    warehouse_client = httpx.AsyncClient(
        transport=httpx.ASGITransport(app=warehouse_app),
        base_url="http://warehouse.test",
    )
    inventory_app = build_inventory_app(warehouse_client, chaos, stats)
    inventory_client = httpx.AsyncClient(
        transport=httpx.ASGITransport(app=inventory_app),
        base_url="http://inventory.test",
        # Bulkhead on the client side too: a small, dedicated pool.
        limits=httpx.Limits(max_connections=BULKHEAD_PERMITS),
    )
    breaker = None
    if use_breaker:
        breaker = CircuitBreaker(
            "inventory",
            BreakerConfig(
                failure_rate_threshold=0.5,
                window_seconds=60.0,
                minimum_calls=4,
                open_duration=30.0,  # >> drill duration: no probes mid-drill
                half_open_max_probes=1,
            ),
            on_event=lambda event, fields: events.append((event, fields)),
        )
    catalog_app = build_catalog_app(
        inventory_client, breaker, asyncio.Semaphore(BULKHEAD_PERMITS)
    )
    return DemoSystem(
        apps={"catalog": catalog_app, "inventory": inventory_app,
              "warehouse": warehouse_app},
        chaos=chaos,
        stats=stats,
        events=events,
    )


async def run_drill(
    requests: int = 30, *, use_breaker: bool, sku: str = "NR-1001"
) -> DrillResult:
    """Fire ``requests`` concurrent catalog calls at a broken inventory."""
    system = build_system(use_breaker=use_breaker)
    system.chaos.always_fail = True  # inventory is hard-down for the drill
    started = time.perf_counter()
    async with httpx.AsyncClient(
        transport=httpx.ASGITransport(app=system.apps["catalog"]),
        base_url="http://catalog.test",
    ) as catalog:
        responses = await asyncio.gather(
            *(catalog.get(f"/products/{sku}") for _ in range(requests))
        )
    elapsed = time.perf_counter() - started
    errors = sum(1 for r in responses if r.status_code >= 500)
    open_rejections = sum(
        1 for r in responses if r.json().get("circuit_open") is True
    )
    return DrillResult(
        requests=requests,
        catalog_errors=errors,
        circuit_open_rejections=open_rejections,
        inventory_calls=system.stats.inventory_calls,
        warehouse_calls=system.stats.warehouse_calls,
        wall_seconds=elapsed,
    )


async def _main() -> None:
    print("Chaos drill: inventory hard-down, 30 concurrent catalog requests")
    for use_breaker in (False, True):
        result = await run_drill(use_breaker=use_breaker)
        label = "WITH breaker " if use_breaker else "WITHOUT breaker"
        print(
            f"{label} | catalog 5xx: {result.catalog_errors:>2}/"
            f"{result.requests} | circuit-open fast-fails: "
            f"{result.circuit_open_rejections:>2} | calls reaching inventory: "
            f"{result.inventory_calls:>2} | wall: {result.wall_seconds:.3f}s"
        )


if __name__ == "__main__":
    asyncio.run(_main())
\end{minted}

Running the drill ({\small\texttt{python demo\_services.py}}) produces the
measured before/after comparison of Table~\ref{tab:ch15-drill}. The
breaker does not make requests succeed --- inventory is genuinely down, and
all thirty catalog calls return 503 either way. What changes is
\emph{where the load goes}: without the breaker, all thirty requests reach
the sick service; with it, four do (the minimum-calls floor plus the
concurrent admissions before the window evaluates) and twenty-six fail fast
and cheap. Those twenty-six unanswered invitations are the headroom in
which inventory recovers.

\begin{table}[h]
  \centering
  \small
  \begin{tabularx}{\textwidth}{l c c X}
    \toprule
    \textbf{Configuration} & \textbf{Calls reaching inventory} & \textbf{Catalog 5xx} & \textbf{Fast-fail rejections} \\
    \midrule
    Without breaker & 30 / 30 & 30 & 0 --- every request pays full failure latency \\
    With breaker & 4 / 30 & 30 & 26 --- rejected in microseconds, dependency spared \\
    \bottomrule
  \end{tabularx}
  \caption{Measured chaos-drill results, Listing~\ref{lst:ch15-demo}
  (deterministic: counts, not wall-clock). Resilience patterns do not
  convert failure into success; they convert expensive failure into cheap
  failure and buy recovery headroom.}
  \label{tab:ch15-drill}
\end{table}

\noindent\textbf{Listing 15.6 --- \texttt{test\_demo\_services.py}:
drill assertions on counts and breaker events.}
\label{lst:ch15-demo-tests}
\begin{minted}{python}
"""Deterministic tests for the three-service chaos drill.

Assertions are on COUNTS (how many requests reached the broken inventory),
never on wall-clock latency. The 30 catalog calls run concurrently via
``asyncio.gather``; the failure mode under test (inventory hard-down,
immediate 500) is fully deterministic.
"""
from __future__ import annotations

import asyncio

import httpx

from demo_services import (
    BULKHEAD_PERMITS,
    build_system,
    run_drill,
)


def test_happy_path_catalog_joins_inventory_and_warehouse() -> None:
    async def main() -> None:
        system = build_system(use_breaker=True)  # chaos OFF
        async with httpx.AsyncClient(
            transport=httpx.ASGITransport(app=system.apps["catalog"]),
            base_url="http://catalog.test",
        ) as client:
            response = await client.get("/products/NR-1001")
        assert response.status_code == 200
        body = response.json()
        assert body["name"] == "Warehouse Picker Bot"
        assert body["availability"] == {"sku": "NR-1001", "available": 5}
        assert system.stats.inventory_calls == 1
        assert system.stats.warehouse_calls == 1

    asyncio.run(main())


def test_without_breaker_every_request_reaches_the_broken_inventory() -> None:
    result = asyncio.run(run_drill(requests=30, use_breaker=False))
    assert result.catalog_errors == 30
    assert result.inventory_calls == 30   # full load hits the sick dependency
    assert result.circuit_open_rejections == 0


def test_with_breaker_most_requests_fail_fast() -> None:
    result = asyncio.run(run_drill(requests=30, use_breaker=True))
    assert result.catalog_errors == 30    # still errors: the dependency IS down
    # ...but almost none touch inventory: breaker opens after 4 failures
    # (minimum_calls, 100% failure rate >= 0.5 threshold).
    assert result.inventory_calls < result.requests
    assert result.circuit_open_rejections == 30 - result.inventory_calls
    assert result.circuit_open_rejections >= 30 - BULKHEAD_PERMITS


def test_breaker_state_transitions_are_emitted_during_drill() -> None:
    system = build_system(use_breaker=True)
    system.chaos.always_fail = True

    async def main() -> None:
        async with httpx.AsyncClient(
            transport=httpx.ASGITransport(app=system.apps["catalog"]),
            base_url="http://catalog.test",
        ) as client:
            await asyncio.gather(*(client.get("/products/NR-1001") for _ in range(30)))

    asyncio.run(main())
    transitions = [f for e, f in system.events if e == "state_change"]
    assert {"previous": "closed", "current": "open"} in [
        {"previous": t["previous"], "current": t["current"]} for t in transitions
    ]
    # Inventory hard-down and breaker open-duration (30s) >> drill duration:
    # exactly one transition, and no half-open probes mid-drill.
    assert len(transitions) == 1
\end{minted}

\subsection{The Transactional Outbox with Relay and Idempotent Consumer}
\label{subsec:ch15-code-outbox}

Listing~\ref{lst:ch15-outbox} implements the outbox pattern end to end
against SQLite. \texttt{OrderService.place\_order} commits the domain row
and the outbox row in one \texttt{with conn} transaction --- atomicity is
the entire point, and the first test proves rollback leaves no orphan
outbox row. \texttt{OutboxRelay.relay\_once} publishes each pending event
and marks it published strictly \emph{after} the publish call returns: the
crash window between those two statements is where at-least-once delivery
comes from, and \texttt{CrashAfterPublish} exists to stand inside it.
\texttt{InventoryConsumer.handle} records the event ID and applies the
stock decrement in one transaction, so a replayed event is absorbed, not
reapplied.

\noindent\textbf{Listing 15.7 --- \texttt{outbox.py}:
atomic domain write + outbox, relay loop, idempotent consumer.}
\label{lst:ch15-outbox}
\begin{minted}{python}
"""Transactional outbox: atomic domain-write + outbox-insert, a relay loop
with at-least-once delivery, and an idempotent consumer.

SQLite-backed, fully in-process. The "broker" is a Python callable, which
keeps the pattern's GUARANTEES visible: the relay marks an event published
only after the publish call returns, so a crash in between replays the
event -- and the consumer must deduplicate (Chapter 10 idempotency).
"""
from __future__ import annotations

import json
import sqlite3
import uuid
from dataclasses import dataclass
from typing import Any, Protocol

PRODUCER_SCHEMA = """
CREATE TABLE IF NOT EXISTS orders (
    order_id  TEXT PRIMARY KEY,
    sku       TEXT NOT NULL,
    quantity  INTEGER NOT NULL CHECK (quantity > 0),
    status    TEXT NOT NULL
);
CREATE TABLE IF NOT EXISTS outbox (
    seq          INTEGER PRIMARY KEY AUTOINCREMENT,
    event_id     TEXT NOT NULL UNIQUE,
    aggregate_id TEXT NOT NULL,
    event_type   TEXT NOT NULL,
    payload      TEXT NOT NULL,
    published    INTEGER NOT NULL DEFAULT 0
);
"""

CONSUMER_SCHEMA = """
CREATE TABLE IF NOT EXISTS processed_events (
    event_id TEXT PRIMARY KEY
);
CREATE TABLE IF NOT EXISTS stock (
    sku      TEXT PRIMARY KEY,
    on_hand  INTEGER NOT NULL
);
"""


class Publisher(Protocol):
    """Anything that can accept one event. A real deployment: Kafka, etc."""

    def publish(self, event: dict[str, Any]) -> None: ...


class InProcessBroker:
    """Deterministic stand-in broker: records every published event."""

    def __init__(self) -> None:
        self.messages: list[dict[str, Any]] = []

    def publish(self, event: dict[str, Any]) -> None:
        self.messages.append(event)


class CrashAfterPublish:
    """Fault-injection publisher: delivers, THEN crashes before the relay
    can mark the event published -- the classic at-least-once window."""

    def __init__(self, inner: Publisher) -> None:
        self._inner = inner

    def publish(self, event: dict[str, Any]) -> None:
        self._inner.publish(event)
        raise RuntimeError("simulated relay crash after publish")


class OrderService:
    """Producer side: domain write and outbox insert in ONE transaction."""

    def __init__(self, conn: sqlite3.Connection) -> None:
        self._conn = conn
        self._conn.executescript(PRODUCER_SCHEMA)

    def place_order(
        self, order_id: str, sku: str, quantity: int, *, event_id: str | None = None
    ) -> str:
        event_id = event_id or uuid.uuid4().hex
        payload = json.dumps({"sku": sku, "quantity": quantity})
        with self._conn:  # atomic: both rows or neither
            self._conn.execute(
                "INSERT INTO orders (order_id, sku, quantity, status)"
                " VALUES (?, ?, ?, 'PLACED')",
                (order_id, sku, quantity),
            )
            self._conn.execute(
                "INSERT INTO outbox (event_id, aggregate_id, event_type, payload)"
                " VALUES (?, ?, 'order.placed', ?)",
                (event_id, order_id, payload),
            )
        return event_id

    def pending_events(self) -> list[dict[str, Any]]:
        rows = self._conn.execute(
            "SELECT event_id, aggregate_id, event_type, payload"
            " FROM outbox WHERE published = 0 ORDER BY seq"
        ).fetchall()
        return [
            {
                "event_id": row[0],
                "aggregate_id": row[1],
                "event_type": row[2],
                "payload": json.loads(row[3]),
            }
            for row in rows
        ]


@dataclass
class OutboxRelay:
    """Polls the outbox, publishes, marks published. At-least-once: the
    mark happens strictly AFTER publish returns."""

    conn: sqlite3.Connection
    publisher: Publisher
    batch_size: int = 100

    def relay_once(self) -> int:
        rows = self.conn.execute(
            "SELECT event_id, aggregate_id, event_type, payload"
            " FROM outbox WHERE published = 0 ORDER BY seq LIMIT ?",
            (self.batch_size,),
        ).fetchall()
        delivered = 0
        for event_id, aggregate_id, event_type, payload in rows:
            event = {
                "event_id": event_id,
                "aggregate_id": aggregate_id,
                "event_type": event_type,
                "payload": json.loads(payload),
            }
            self.publisher.publish(event)  # crash here -> replay later
            with self.conn:
                self.conn.execute(
                    "UPDATE outbox SET published = 1 WHERE event_id = ?",
                    (event_id,),
                )
            delivered += 1
        return delivered


class InventoryConsumer:
    """Idempotent consumer: effect applied at most once, tracked in the
    same transaction as the dedup record."""

    def __init__(self, conn: sqlite3.Connection) -> None:
        self._conn = conn
        self._conn.executescript(CONSUMER_SCHEMA)

    def seed_stock(self, sku: str, on_hand: int) -> None:
        with self._conn:
            self._conn.execute(
                "INSERT INTO stock (sku, on_hand) VALUES (?, ?)"
                " ON CONFLICT(sku) DO UPDATE SET on_hand = excluded.on_hand",
                (sku, on_hand),
            )

    def on_hand(self, sku: str) -> int:
        row = self._conn.execute(
            "SELECT on_hand FROM stock WHERE sku = ?", (sku,)
        ).fetchone()
        return row[0] if row else 0

    def handle(self, event: dict[str, Any]) -> bool:
        """Apply the event; return False for duplicates (already processed)."""
        if event["event_type"] != "order.placed":
            raise ValueError(f"unsupported event type {event['event_type']!r}")
        with self._conn:
            cursor = self._conn.execute(
                "INSERT OR IGNORE INTO processed_events (event_id) VALUES (?)",
                (event["event_id"],),
            )
            if cursor.rowcount == 0:
                return False  # duplicate delivery: deduped
            self._conn.execute(
                "UPDATE stock SET on_hand = on_hand - ? WHERE sku = ?",
                (event["payload"]["quantity"], event["payload"]["sku"]),
            )
            return True
\end{minted}

\noindent\textbf{Listing 15.8 --- \texttt{test\_outbox.py}:
atomicity, crash-mid-publish replay, and dedup proofs.}
\label{lst:ch15-outbox-tests}
\begin{minted}{python}
"""Tests for the transactional outbox, relay crash, and consumer dedup."""
from __future__ import annotations

import sqlite3

import pytest

from outbox import (
    CrashAfterPublish,
    InProcessBroker,
    InventoryConsumer,
    OrderService,
    OutboxRelay,
)


@pytest.fixture()
def producer_conn() -> sqlite3.Connection:
    return sqlite3.connect(":memory:")


@pytest.fixture()
def consumer() -> InventoryConsumer:
    consumer = InventoryConsumer(sqlite3.connect(":memory:"))
    consumer.seed_stock("NR-1001", 100)
    return consumer


def test_domain_write_and_outbox_insert_are_atomic(producer_conn) -> None:
    service = OrderService(producer_conn)
    service.place_order("ORD-1", "NR-1001", 3, event_id="evt-1")
    # Second insert with the same order_id violates the PK; the outbox row
    # from the failed transaction must be rolled back too.
    with pytest.raises(sqlite3.IntegrityError):
        service.place_order("ORD-1", "NR-1001", 5, event_id="evt-2")
    pending = service.pending_events()
    assert [e["event_id"] for e in pending] == ["evt-1"]


def test_relay_delivers_pending_events_and_marks_them(producer_conn) -> None:
    service = OrderService(producer_conn)
    service.place_order("ORD-1", "NR-1001", 3, event_id="evt-1")
    service.place_order("ORD-2", "NR-2002", 1, event_id="evt-2")
    broker = InProcessBroker()
    relay = OutboxRelay(producer_conn, broker)
    assert relay.relay_once() == 2
    assert [m["event_id"] for m in broker.messages] == ["evt-1", "evt-2"]
    assert relay.relay_once() == 0  # nothing left pending


def test_relay_crash_mid_publish_replays_event(producer_conn) -> None:
    service = OrderService(producer_conn)
    service.place_order("ORD-1", "NR-1001", 3, event_id="evt-1")
    broker = InProcessBroker()

    crashing = OutboxRelay(producer_conn, CrashAfterPublish(broker))
    with pytest.raises(RuntimeError, match="simulated relay crash"):
        crashing.relay_once()
    assert len(broker.messages) == 1       # delivered...
    assert len(service.pending_events()) == 1  # ...but NOT marked published

    recovered = OutboxRelay(producer_conn, broker)
    assert recovered.relay_once() == 1     # at-least-once: replayed
    assert len(broker.messages) == 2       # duplicate on the wire
    assert len(service.pending_events()) == 0


def test_idempotent_consumer_dedups_replayed_events(consumer) -> None:
    event = {
        "event_id": "evt-1",
        "aggregate_id": "ORD-1",
        "event_type": "order.placed",
        "payload": {"sku": "NR-1001", "quantity": 3},
    }
    assert consumer.handle(event) is True
    assert consumer.on_hand("NR-1001") == 97
    assert consumer.handle(event) is False  # duplicate: no double decrement
    assert consumer.handle(event) is False
    assert consumer.on_hand("NR-1001") == 97


def test_end_to_end_order_to_stock_under_crash_and_replay(
    producer_conn, consumer
) -> None:
    service = OrderService(producer_conn)
    service.place_order("ORD-1", "NR-1001", 3, event_id="evt-1")
    service.place_order("ORD-2", "NR-1001", 2, event_id="evt-2")
    broker = InProcessBroker()

    crashing = OutboxRelay(producer_conn, CrashAfterPublish(broker))
    with pytest.raises(RuntimeError):
        crashing.relay_once()  # evt-1 delivered, crash before mark

    OutboxRelay(producer_conn, broker).relay_once()  # replays evt-1, adds evt-2
    assert [m["event_id"] for m in broker.messages] == ["evt-1", "evt-1", "evt-2"]

    applied = [consumer.handle(m) for m in broker.messages]
    assert applied == [True, False, True]
    assert consumer.on_hand("NR-1001") == 95  # 3 + 2 deducted exactly once
\end{minted}

The end-to-end test deserves a slow reading because it \emph{is} the
pattern's guarantee: two orders are placed; the relay delivers
\texttt{evt-1} and crashes before marking it; the recovered relay replays
\texttt{evt-1} and delivers \texttt{evt-2}, so the wire carries
\texttt{[evt-1, evt-1, evt-2]}; the consumer applies
\texttt{[True, False, True]}; and stock ends at exactly
$100 - 3 - 2 = 95$. At-least-once transport plus idempotent consumption
equals effectively-once processing --- no two-phase commit involved.

\subsection{Saga Orchestrator with Compensating Actions}
\label{subsec:ch15-code-saga}

Listing~\ref{lst:ch15-saga} implements the orchestrated saga of
Definition~\ref{def:ch15-saga} for Northwind Robotics order fulfillment:
reserve inventory, charge payment, schedule shipment, each with its
compensation (release, refund, cancel). The orchestrator is deliberately
small --- the value is in the explicitness: steps in order, compensations
in reverse, every transition emitted to the event hook, and a
\textsc{compensation\_failed} status that is \emph{surfaced}, never
swallowed, because a failed compensation means the world is in a state only
an operator can reconcile. The in-memory domain services take a
\texttt{fail\_on} set so tests can break any operation by name.

\noindent\textbf{Listing 15.9 --- \texttt{saga.py}:
orchestrated order-fulfillment saga with compensations and failure
injection.}
\label{lst:ch15-saga}
\begin{minted}{python}
"""Saga orchestrator for Northwind Robotics order fulfillment.

Flow: reserve inventory -> charge payment -> schedule shipment.
Each step declares a compensating action; on failure the orchestrator runs
compensations for completed steps in REVERSE order and reports the saga as
COMPENSATED (not "rolled back" -- the real world already happened; we
compensate, we do not time-travel).

Failure injection: each domain service takes ``fail_on`` -- a set of
operation names that raise on invocation.
"""
from __future__ import annotations

from collections.abc import Callable
from dataclasses import dataclass, field
from enum import Enum
from typing import Any


class SagaStatus(str, Enum):
    COMPLETED = "completed"
    COMPENSATED = "compensated"
    COMPENSATION_FAILED = "compensation_failed"


class CompensationError(RuntimeError):
    """A compensating action itself failed; requires operator attention."""


@dataclass(frozen=True, slots=True)
class SagaStep:
    name: str
    action: Callable[[dict[str, Any]], None]
    compensate: Callable[[dict[str, Any]], None] | None = None


@dataclass(frozen=True, slots=True)
class SagaResult:
    status: SagaStatus
    failed_step: str | None
    compensated: tuple[str, ...]
    context: dict[str, Any]


class SagaOrchestrator:
    """Executes steps in order; compensates in reverse on failure."""

    def __init__(
        self,
        steps: list[SagaStep],
        on_event: Callable[[str, dict[str, Any]], None] | None = None,
    ) -> None:
        if not steps:
            raise ValueError("a saga needs at least one step")
        self._steps = steps
        self._on_event = on_event

    def _emit(self, event: str, **fields: Any) -> None:
        if self._on_event is not None:
            self._on_event(event, fields)

    def run(self, context: dict[str, Any] | None = None) -> SagaResult:
        context = dict(context or {})
        completed: list[SagaStep] = []
        for step in self._steps:
            self._emit("step_started", step=step.name)
            try:
                step.action(context)
            except Exception as exc:
                self._emit("step_failed", step=step.name, error=type(exc).__name__)
                return self._compensate(completed, failed=step.name, context=context)
            completed.append(step)
            self._emit("step_completed", step=step.name)
        return SagaResult(
            status=SagaStatus.COMPLETED,
            failed_step=None,
            compensated=(),
            context=context,
        )

    def _compensate(
        self, completed: list[SagaStep], *, failed: str, context: dict[str, Any]
    ) -> SagaResult:
        compensated: list[str] = []
        try:
            for step in reversed(completed):
                if step.compensate is None:
                    continue
                self._emit("compensation_started", step=step.name)
                step.compensate(context)
                compensated.append(step.name)
                self._emit("compensation_completed", step=step.name)
        except Exception as exc:
            self._emit(
                "compensation_failed", step=step.name, error=type(exc).__name__
            )
            return SagaResult(
                status=SagaStatus.COMPENSATION_FAILED,
                failed_step=failed,
                compensated=tuple(compensated),
                context=context,
            )
        return SagaResult(
            status=SagaStatus.COMPENSATED,
            failed_step=failed,
            compensated=tuple(compensated),
            context=context,
        )


# --------------------------------------------------------------------------
# Northwind Robotics domain services (in-memory, deterministic)
# --------------------------------------------------------------------------
class InventoryService:
    def __init__(self, stock: dict[str, int], fail_on: set[str] | None = None) -> None:
        self.stock = dict(stock)
        self.reserved: dict[str, int] = {}
        self.fail_on = fail_on or set()
        self.calls: list[str] = []

    def reserve(self, sku: str, quantity: int) -> None:
        self.calls.append("reserve")
        if "reserve" in self.fail_on:
            raise ConnectionError("inventory service unavailable")
        if self.stock.get(sku, 0) < quantity:
            raise ValueError(f"insufficient stock for {sku}")
        self.stock[sku] -= quantity
        self.reserved[sku] = self.reserved.get(sku, 0) + quantity

    def release(self, sku: str, quantity: int) -> None:
        self.calls.append("release")
        if "release" in self.fail_on:
            raise ConnectionError("inventory service unavailable")
        self.reserved[sku] = self.reserved.get(sku, 0) - quantity
        self.stock[sku] = self.stock.get(sku, 0) + quantity


class PaymentService:
    def __init__(self, fail_on: set[str] | None = None) -> None:
        self.charges: list[tuple[str, float]] = []
        self.refunds: list[tuple[str, float]] = []
        self.fail_on = fail_on or set()
        self.calls: list[str] = []

    def charge(self, order_id: str, amount: float) -> None:
        self.calls.append("charge")
        if "charge" in self.fail_on:
            raise ConnectionError("payment gateway timeout")
        self.charges.append((order_id, amount))

    def refund(self, order_id: str, amount: float) -> None:
        self.calls.append("refund")
        if "refund" in self.fail_on:
            raise ConnectionError("payment gateway timeout")
        self.refunds.append((order_id, amount))


class ShipmentService:
    def __init__(self, fail_on: set[str] | None = None) -> None:
        self.shipments: list[str] = []
        self.cancellations: list[str] = []
        self.fail_on = fail_on or set()
        self.calls: list[str] = []

    def schedule(self, order_id: str) -> None:
        self.calls.append("schedule")
        if "schedule" in self.fail_on:
            raise ConnectionError("carrier API error")
        self.shipments.append(order_id)

    def cancel(self, order_id: str) -> None:
        self.calls.append("cancel")
        if "cancel" in self.fail_on:
            raise ConnectionError("carrier API error")
        self.cancellations.append(order_id)


def build_order_fulfillment_saga(
    inventory: InventoryService,
    payments: PaymentService,
    shipments: ShipmentService,
    on_event: Callable[[str, dict[str, Any]], None] | None = None,
) -> SagaOrchestrator:
    """The canonical three-step order saga with compensations."""
    steps = [
        SagaStep(
            name="reserve_inventory",
            action=lambda ctx: inventory.reserve(ctx["sku"], ctx["quantity"]),
            compensate=lambda ctx: inventory.release(ctx["sku"], ctx["quantity"]),
        ),
        SagaStep(
            name="charge_payment",
            action=lambda ctx: payments.charge(ctx["order_id"], ctx["amount"]),
            compensate=lambda ctx: payments.refund(ctx["order_id"], ctx["amount"]),
        ),
        SagaStep(
            name="schedule_shipment",
            action=lambda ctx: shipments.schedule(ctx["order_id"]),
            compensate=lambda ctx: shipments.cancel(ctx["order_id"]),
        ),
    ]
    return SagaOrchestrator(steps, on_event=on_event)
\end{minted}

\noindent\textbf{Listing 15.10 --- \texttt{test\_saga.py}:
failure-injection proofs, including reverse-order compensation and
surfaced compensation failure.}
\label{lst:ch15-saga-tests}
\begin{minted}{python}
"""Failure-injection tests for the saga orchestrator."""
from __future__ import annotations

from saga import (
    InventoryService,
    PaymentService,
    SagaOrchestrator,
    SagaStatus,
    SagaStep,
    ShipmentService,
    build_order_fulfillment_saga,
)

ORDER = {"order_id": "ORD-7", "sku": "NR-1001", "quantity": 3, "amount": 184.50}


def build(fail_on: dict[str, set[str]] | None = None):
    fail_on = fail_on or {}
    inventory = InventoryService({"NR-1001": 10}, fail_on.get("inventory"))
    payments = PaymentService(fail_on.get("payments"))
    shipments = ShipmentService(fail_on.get("shipments"))
    events: list[tuple[str, dict]] = []
    saga = build_order_fulfillment_saga(
        inventory, payments, shipments,
        on_event=lambda event, fields: events.append((event, fields)),
    )
    return saga, inventory, payments, shipments, events


def test_happy_path_completes_all_steps_in_order() -> None:
    saga, inventory, payments, shipments, events = build()
    result = saga.run(ORDER)
    assert result.status is SagaStatus.COMPLETED
    assert result.failed_step is None
    assert result.compensated == ()
    assert inventory.stock["NR-1001"] == 7
    assert inventory.reserved["NR-1001"] == 3
    assert payments.charges == [("ORD-7", 184.50)]
    assert shipments.shipments == ["ORD-7"]
    started = [f["step"] for e, f in events if e == "step_started"]
    assert started == ["reserve_inventory", "charge_payment", "schedule_shipment"]


def test_payment_failure_compensates_inventory_in_reverse() -> None:
    saga, inventory, payments, shipments, _ = build(
        {"payments": {"charge"}}
    )
    result = saga.run(ORDER)
    assert result.status is SagaStatus.COMPENSATED
    assert result.failed_step == "charge_payment"
    assert result.compensated == ("reserve_inventory",)
    assert inventory.stock["NR-1001"] == 10   # reservation released
    assert inventory.reserved["NR-1001"] == 0
    assert payments.charges == []             # charge never landed
    assert shipments.shipments == []          # never reached


def test_shipment_failure_compensates_payment_then_inventory() -> None:
    saga, inventory, payments, shipments, _ = build(
        {"shipments": {"schedule"}}
    )
    result = saga.run(ORDER)
    assert result.status is SagaStatus.COMPENSATED
    assert result.failed_step == "schedule_shipment"
    # Reverse order of completion: refund BEFORE release.
    assert result.compensated == ("charge_payment", "reserve_inventory")
    assert payments.charges == [("ORD-7", 184.50)]   # happened, then refunded
    assert payments.refunds == [("ORD-7", 184.50)]
    assert inventory.stock["NR-1001"] == 10
    assert shipments.shipments == []


def test_first_step_failure_needs_no_compensation() -> None:
    saga, inventory, payments, shipments, _ = build(
        {"inventory": {"reserve"}}
    )
    result = saga.run(ORDER)
    assert result.status is SagaStatus.COMPENSATED
    assert result.failed_step == "reserve_inventory"
    assert result.compensated == ()
    assert payments.charges == []
    assert shipments.shipments == []


def test_compensation_failure_is_surfaced_not_swallowed() -> None:
    saga, inventory, payments, shipments, events = build(
        {"payments": {"charge", "refund"}}
    )
    # charge fails -> saga compensates reserve (refund is never reached
    # because charge never completed). To exercise a FAILING compensation,
    # fail schedule + refund.
    saga, inventory, payments, shipments, events = build(
        {"shipments": {"schedule"}, "payments": {"refund"}}
    )
    result = saga.run(ORDER)
    assert result.status is SagaStatus.COMPENSATION_FAILED
    assert result.failed_step == "schedule_shipment"
    # Refund failed, so inventory release never ran: partial compensation
    # is explicit in the result for an operator to reconcile.
    assert result.compensated == ()
    assert any(e == "compensation_failed" for e, _ in events)
    assert inventory.stock["NR-1001"] == 7  # still reserved: needs a human


def test_orchestrator_rejects_empty_step_list() -> None:
    import pytest

    with pytest.raises(ValueError):
        SagaOrchestrator([])


def test_saga_steps_without_compensation_are_skipped() -> None:
    log: list[str] = []

    def fail(_ctx) -> None:
        raise ConnectionError("down")

    saga = SagaOrchestrator(
        [
            SagaStep("fire_and_forget_email", lambda ctx: log.append("email")),
            SagaStep("charge", fail, compensate=lambda ctx: log.append("refund")),
        ]
    )
    result = saga.run({})
    assert result.status is SagaStatus.COMPENSATED
    assert result.compensated == ()  # email step has no compensation; charge never completed
    assert log == ["email"]
\end{minted}

The tests cover the state space that matters: the happy path's step
ordering; payment failure compensated by exactly one release; shipment
failure compensated refund-\emph{then}-release (reverse order, asserted on
the tuple, not on vibes); first-step failure requiring no compensation at
all; and the uncomfortable one --- a compensation that itself fails,
leaving stock reserved and the saga in \textsc{compensation\_failed} with
the partial compensation list intact for reconciliation. That last test is
the difference between a saga implementation and a saga demo.

%% ----------------------------------------------------------------------------
\section{Common Pitfalls \& Anti-Patterns}
\label{sec:ch15-pitfalls}

Each anti-pattern below has been the proximal cause of a real outage class.
Format: what it looks like, why it fails, the fix.

\subsection*{P15.1 --- The Distributed Monolith via Synchronous Chains}
\emph{Looks like:} eight services, and rendering one page calls all of
them, in sequence, synchronously. \emph{Why it fails:} availability
multiplies downward --- a chain of $n$ services each at $99.9\%$
availability offers the caller at most $0.999^n \approx 99.2\%$ for
$n = 8$, before latency tails and retry amplification are counted. You have
deployed a monolith with network faults inside its function calls.
\emph{Fix:} shorten synchronous chains ruthlessly; move non-critical work
behind events (Section~\ref{subsec:ch15-taxonomy}); replicate read data
locally instead of fetching it per request; ask of every hop whether the
response is needed \emph{before} the caller's deadline or merely
\emph{eventually}.

\subsection*{P15.2 --- Retries without Budget, Jitter, or Idempotency}
\emph{Looks like:} \texttt{for attempt in range(5): try: ...} at every
layer, fixed 100\,ms backoff, applied to every endpoint including
non-idempotent POSTs. \emph{Why it fails:}
Theorem~\ref{thm:ch15-retry-amplification} --- you have built a $39\times$
load multiplier that engages precisely when the system is least able to
absorb load, synchronized across clients by the shared fixed backoff, and
applied to operations that must not run twice. \emph{Fix:} the
Listing~\ref{lst:ch15-retry-budget} discipline: retry only idempotent
operations or those carrying idempotency keys
(Chapter~\ref{ch:idempotency}), spend from a process-wide token bucket,
jitter fully, respect deadlines, honor \texttt{Retry-After}.

\subsection*{P15.3 --- One Breaker Shared across Dependencies}
\emph{Looks like:} a single global \texttt{CircuitBreaker("downstream")}
wrapping calls to inventory, payments, and recommendations alike.
\emph{Why it fails:} the breaker's state is a hypothesis about one
resource; a sick recommendations service now opens the circuit for payments
--- converting a partial degradation into a total one. \emph{Fix:} one
breaker per dependency per caller (Listing~\ref{lst:ch15-breaker} enforces
this by requiring a name), each with thresholds tuned to that dependency's
traffic volume and SLO.

\subsection*{P15.4 --- A Timeout Longer Than the Caller's}
\emph{Looks like:} the edge times out at $2\,\mathrm{s}$; catalog calls
inventory with a $5\,\mathrm{s}$ timeout ``to be safe.'' \emph{Why it
fails:} after the edge abandons the request, catalog keeps working for
three more seconds --- burning bulkhead permits, connections, and inventory
capacity on work nobody will read. Under load this is stage 2--3 of
Figure~\ref{fig:ch15-cascade}. \emph{Fix:}
Definition~\ref{def:ch15-deadline}: budgets shrink strictly toward the
leaves; propagate deadlines (gRPC does this on the wire~\cite{grpcdocs});
verify with a drill.

\subsection*{P15.5 --- Ignoring Bulkheads}
\emph{Looks like:} all endpoints share one connection pool, one thread
pool, one event loop's worth of unbounded concurrency into every
dependency. \emph{Why it fails:} the slowest dependency consumes the shared
pool, and now healthy dependencies fail too --- the failure domain is the
whole process. \emph{Fix:} per-dependency clients with their own
\texttt{Limits}, semaphores per dependency
(Listing~\ref{lst:ch15-demo}), and reserve capacity for health checks and
metrics so the control surfaces keep working during the incident.

\subsection*{P15.6 --- Two-Phase Commit over HTTP}
\emph{Looks like:} \texttt{POST /prepare} on three services, then
\texttt{POST /commit} on all, orchestrated over stateless HTTP.
\emph{Why it fails:} HTTP gives you no way to hold a prepare lock across
requests, no coordinator-crash protocol, and no answer when the commit
response is lost; you have re-implemented 2PC's blocking, availability-
coupling core with none of its correctness machinery. \emph{Fix:}
Section~\ref{subsec:ch15-transactions}: sagas with compensations, the
outbox for event emission, idempotent consumers. Accept eventual
consistency as the design input it is.

\subsection*{P15.7 --- Assuming the Mesh Fixes Application Semantics}
\emph{Looks like:} ``we have Istio, so we can delete the retry logic, the
idempotency keys, and the saga code.'' \emph{Why it fails:} the mesh sees
packets, not meaning. It will happily retry your non-idempotent charge
endpoint three times (now at infrastructure level, where your code review
can't see it), and it cannot express ``compensate reserve-inventory with
release.'' \emph{Fix:} treat the mesh as the transport-hygiene layer
(mTLS, telemetry, per-hop timeouts, LB) and keep application-level
correctness --- idempotency, sagas, deadline budgets --- in the
application, where the semantics live.

\subsection*{P15.8 --- Breakers and Budgets without Metrics}
\emph{Looks like:} a breaker that opens in production, and the first
anyone learns of it is the 503 spike in a customer screenshot.
\emph{Why it fails:} every pattern in this chapter is a control system; a
control system you cannot observe is a superstition. \emph{Fix:} wire the
\texttt{on\_event} hooks of Listings~\ref{lst:ch15-breaker}
and~\ref{lst:ch15-retry-budget} to your metrics pipeline; alert on
\texttt{state\_change} to open, on sustained \texttt{budget\_exhausted},
and on probe-failure loops (open $\to$ half-open $\to$ open flapping).

\begin{warningbox}
The composite anti-pattern is \emph{pattern theater}: breakers, retries,
and timeouts present in code review but untested under failure. If you have
never watched your breaker open in a drill, you do not have a breaker; you
have a hope. Section~\ref{sec:ch15-lab} is the remedy.
\end{warningbox}

%% ----------------------------------------------------------------------------
\section{Hands-On Production Lab \& Real-World Case Study}
\label{sec:ch15-lab}

\subsection*{Lab: Run the Chaos Drill, Watch the Breaker, Tune the Thresholds}

Work in the chapter code directory. Everything runs offline.

\begin{enumerate}
  \item \textbf{Baseline.} Create and populate the environment, then run
        the full suite and the drill:
\begin{minted}{text}
python -m venv .venv
.\.venv\Scripts\Activate.ps1
pip install -r requirements.txt
pytest -q
python demo_services.py
\end{minted}
        Confirm the drill reproduces Table~\ref{tab:ch15-drill}: 30 calls
        reach the broken inventory without the breaker, 4 with it.
  \item \textbf{Observe state transitions.} In \texttt{demo\_services.py},
        the breaker is constructed with an \texttt{on\_event} hook feeding
        \texttt{DemoSystem.events}. Extend \texttt{\_main} to print every
        \texttt{state\_change} and \texttt{rejected} event with its
        payload, and re-run. You should see exactly one
        closed~$\to$~open transition and no half-open probes, because the
        drill is shorter than \texttt{open\_duration}.
  \item \textbf{Watch a probe.} Shorten \texttt{open\_duration} to
        \texttt{0.0}-ish (use \texttt{0.05}) and insert
        \texttt{await asyncio.sleep(0.1)} into the drill loop after the
        first batch; flip \texttt{chaos.always\_fail} off mid-drill. Now a
        half-open probe succeeds and the breaker closes --- observe the
        open~$\to$~half\_open~$\to$~closed sequence in the events, and
        catalog returning 200s again.
  \item \textbf{Tune thresholds.} Sweep \texttt{minimum\_calls} over
        $\{2, 4, 8\}$ and \texttt{failure\_rate\_threshold} over
        $\{0.25, 0.5, 1.0\}$. Record calls-reaching-inventory for each
        combination. Lower floors open sooner (better protection, more
        false positives under transient single-failure noise); higher
        thresholds tolerate flaky-but-alive dependencies.
  \item \textbf{Break the budget rule.} Set \texttt{CATALOG\_TIMEOUT}
        larger than the caller's deadline and re-run the drill with
        latency chaos (\texttt{chaos.latency\_seconds = 0.4}). Count
        abandoned work. Restore the rule.
\end{enumerate}

\subsection*{War Story: The 200 ms Blip That Took Down Three Services}

Northwind Robotics, a Tuesday. The warehouse service's storage layer
hiccupped: $p99$ latency jumped from $40\,\mathrm{ms}$ to
$200\,\mathrm{ms}$ for ninety seconds --- a blip, by any dashboard's
standards. Here is what the composition did with it.

Inventory's calls to warehouse, budgeted at $100\,\mathrm{ms}$, now timed
out; inventory's client retried twice with a fixed $50\,\mathrm{ms}$
backoff. Catalog called inventory with retries of its own. Within forty
seconds the warehouse was receiving roughly $27\times$ its normal request
rate --- Theorem~\ref{thm:ch15-retry-amplification} with $a = 3$, $n = 3$,
and the feedback loop of Figure~\ref{fig:ch15-cascade} closing: each retry
wave deepened the warehouse's queue, which raised latency further, which
converted more slow successes into timeout ``failures,'' which justified
more retries. The warehouse, which had been \emph{slow}, was now genuinely
down; inventory's shared connection pool exhausted waiting on it, so
inventory stopped serving \emph{all} requests, including the ones that
never touched warehouse; catalog followed. Three services down, customer-
visible, from a 200\,ms blip. The blip had healed on its own within ninety
seconds; the outage took forty minutes, because each restart walked
straight back into $27\times$ load.

The postmortem fixes read as this chapter's table of contents: per-hop
deadline budgets that shrink toward the leaves; full jitter on every
backoff; a process-wide retry budget of ten percent of request volume; one
circuit breaker per dependency with a minimum-call floor; a dedicated
warehouse client pool so inventory's other endpoints kept their own
capacity; and a standing chaos drill so the next 200\,ms blip is a line in
a metrics dashboard instead of an incident review. The blips did not stop.
The outages did.

%% ----------------------------------------------------------------------------
\section{Self-Assessment Exercises \& Deep-Dive Questions}
\label{sec:ch15-exercises}

\begin{enumerate}
  \item \textbf{Deadline budgets.} A public API promises $p99 < 1.5
        \,\mathrm{s}$. The chain is edge $\to$ orders $\to$ payments $\to$
        ledger, and ledger's measured $p99.9$ is $300\,\mathrm{ms}$.
        Propose per-hop budgets satisfying Definition~\ref{def:ch15-deadline}
        with headroom at each layer, and state what you would measure to
        validate them.
  \item \textbf{Amplification audit.} Your platform's call graph has four
        hops; each hop currently retries up to twice with no budget.
        Compute the worst-case downstream load multiplier. Now impose a
        one-token retry budget per process at the two middle hops: recompute.
        Which hop's budget matters most, and why?
  \item \textbf{Breaker parameters.} A dependency serves
        $2{,}000\,\mathrm{req/s}$ to your fleet at $99.95\%$ success, and
        its team asks you not to exceed $20\,\mathrm{req/s}$ during its
        recoveries. Design $(\theta, W, N, d, p, s)$ from
        Definition~\ref{def:ch15-breaker} and justify each number.
  \item \textbf{Extend the breaker.} Add a \emph{slow-call} dimension to
        Listing~\ref{lst:ch15-breaker}: calls slower than a threshold count
        as failures even when they succeed. Prove with a manual-clock test
        that five slow successes open the breaker.
  \item \textbf{Hedge, honestly.} Extend the demo's catalog $\to$ inventory
        call to hedge after the measured $p95$. Instrument the extra-call
        rate. At what injected $q = \Pr[T > h]$ does your measured hedge
        rate match $1 + q$? What must be true of the endpoint for hedging
        to be safe?
  \item \textbf{Outbox ordering.} The relay of
        Listing~\ref{lst:ch15-outbox} publishes in \texttt{seq} order, but
        a crash mid-batch plus a retry can interleave replays with new
        events. Specify (and test) a per-aggregate ordering guarantee:
        what changes in the outbox schema, the relay, and the consumer?
  \item \textbf{Choreograph the saga.} Re-express
        Listing~\ref{lst:ch15-saga}'s flow as choreography over the
        in-process broker: inventory emits \texttt{inventory.reserved},
        payments consumes it, and so on. Where does the compensation graph
        now live? Write the test that catches a missing compensation
        subscription.
  \item \textbf{Mesh threat model.} Your mesh provides mTLS everywhere.
        Enumerate three application-level failure or misuse modes that
        mTLS does nothing for, and name the chapter pattern that addresses
        each.
\end{enumerate}

%% ----------------------------------------------------------------------------
\section{Interview Q\&A Vault}
\label{sec:ch15-interview}

\subsection*{Junior (Q1--Q10)}
\begin{enumerate}[label=\textbf{Q\arabic*.}, leftmargin=1.6cm]
  \item \textbf{Synchronous versus asynchronous communication between
        services --- what's the difference and when do you pick each?}
        Synchronous (REST, gRPC): the caller blocks for a response; simple
        mental model, tight temporal coupling, additive latency. Asynchronous
        (queues, streams): the caller hands work to a broker; low temporal
        coupling, natural fan-out and load smoothing, at the price of
        duplicates, lag, and harder status queries. Pick synchronous for
        user-facing queries that need an answer now; asynchronous for
        side effects, fan-out, and anything whose consumer can be down.
  \item \textbf{What is a timeout, and why is ``no timeout'' never the
        right default?}
        A timeout is an upper bound on how long a call may occupy the
        caller's resources. Without one, a slow dependency holds your
        threads, connections, and memory indefinitely --- slow is worse
        than down, because down gives you an error to react to and slow
        gives you a resource leak.
  \item \textbf{What are the three states of a circuit breaker?}
        \textsc{Closed}: calls flow and failures are counted.
        \textsc{Open}: calls are rejected immediately without touching the
        dependency. \textsc{Half\_open}: a limited number of probe calls
        are admitted to test recovery; success closes, failure reopens.
  \item \textbf{Why does a circuit breaker need a minimum number of calls
        before it can open?}
        Statistics. Two failures out of two overnight calls must not trip
        a circuit that carries thousands of calls per minute at noon. The
        minimum-call floor $N$ ensures the failure rate is evaluated only
        once the window contains enough evidence to mean something.
  \item \textbf{What is a bulkhead, in software?}
        A partition of resources --- connection pools, thread pools,
        concurrency permits --- so that one failing dependency cannot
        consume the capacity everything else needs. Named for ship
        compartments: a breach floods one compartment, not the hull.
  \item \textbf{What is the difference between orchestration and
        choreography?}
        Orchestration: one coordinator explicitly drives each step and owns
        workflow state. Choreography: services react to events and emit
        their own; the workflow is emergent. Orchestration centralizes
        failure handling and status queries; choreography removes the
        coordinator as a bottleneck and single point of failure.
  \item \textbf{What does it mean for a consumer to be idempotent, and why
        does at-least-once delivery require it?}
        An idempotent consumer processes the same event many times with the
        effect of processing it once --- typically by recording processed
        event IDs in the same transaction as the state change. At-least-once
        delivery guarantees replays; without dedup, replays double-apply
        effects (charging twice, decrementing stock twice).
  \item \textbf{What is service discovery?}
        The mechanism by which callers learn the current network locations
        of a logical service's instances --- via DNS, a registry with
        health checking, or a platform API --- replacing spatially coupled
        hardcoded addresses that break under autoscaling and rescheduling.
  \item \textbf{Name three of the fallacies of distributed computing and a
        consequence of believing each.}
        ``The network is reliable'' $\to$ you cannot distinguish
        unprocessed from unacknowledged requests, so retries need
        idempotency. ``Latency is zero'' $\to$ tail latency, not the mean,
        must set your timeout budgets. ``Topology doesn't change'' $\to$
        hardcoded addresses become outages the first time autoscalers run.
  \item \textbf{What is a compensating action?}
        The semantic undo of a saga step: release the reservation, refund
        the charge, cancel the shipment. It is not a database rollback ---
        the original step really happened and the world saw it; the
        compensation applies a forward action that restores the business
        invariant.
\end{enumerate}

\subsection*{Mid (Q11--Q20)}
\begin{enumerate}[label=\textbf{Q\arabic*.}, leftmargin=1.6cm, start=11]
  \item \textbf{Derive the retry amplification of an $n$-hop chain where
        each hop makes up to $a$ attempts.}
        Let $L_k$ be calls arriving at hop $k$; $L_0 = 1$ (the client
        request) and $L_{k+1} = a \cdot L_k$, so $L_k = a^k$ and the total
        downstream load is $S(n,a) = \sum_{k=1}^{n} a^k =
        a(a^n-1)/(a-1)$. With $a = 3$, $n = 3$: $3 + 9 + 27 = 39$ requests
        per client call --- $39\times$ amplification, $27\times$ at the
        leaf (Theorem~\ref{thm:ch15-retry-amplification}).
  \item \textbf{How do retry budgets cap amplification?}
        By making retries spend from a process-wide token bucket where only
        retries (never first attempts) consume tokens. Attempts per logical
        request are then bounded by $1 + B$ regardless of configured
        \texttt{max\_retries}, and sustained retry rate by the refill rate
        (Proposition~\ref{prop:ch15-budget-cap}).
        Listing~\ref{lst:ch15-retry-budget-tests} proves it: five permitted
        retries, two tokens, three requests total.
  \item \textbf{Why jitter, and why \emph{full} jitter?}
        Failures correlate clients in time; identical deterministic
        backoffs re-correlate them into waves that re-collide with the
        recovering service. Full jitter draws each sleep uniformly from
        $[0, \min(c, b \cdot 2^{k})]$, decorrelating the herd while keeping
        the exponential envelope --- at the cost of occasionally short
        delays, which is what the budget is for.
  \item \textbf{Explain the deadline budget rule for a call chain.}
        Every hop's timeout must be smaller than its caller's remaining
        budget, with headroom: $T_i > \sum_{j > i} T_j$. Violation means
        the caller gives up while the callee keeps working --- pure waste,
        and under load, pool exhaustion. gRPC propagates an absolute
        deadline in metadata so each hop derives its remaining budget
        instead of re-guessing~\cite{grpcdocs}.
  \item \textbf{Design breaker thresholds for a dependency with a 99.9\%
        success SLO serving you 500 req/s.}
        Work backward from the SLO and volume. Baseline failure rate is
        $0.1\%$; set $\theta = 0.5$ so the breaker opens only when failure
        is unambiguously systemic, not at SLO noise. Window $W = 10$ s
        covers $\approx 5{,}000$ calls; set $N = 20$ (a floor far below
        normal volume so statistics are meaningful, far above noise). Open
        duration $d = 5$--$30$ s, tuned to the dependency's observed
        recovery time. Probes: $p = 1$--$3$ concurrent, $s = 3$ consecutive
        successes --- recovery is when probes succeed repeatedly, not once.
        Validate every number with the chaos drill, not with hope.
  \item \textbf{What is the tail-amplification trade-off of hedged
        requests?}
        Hedging fires a second identical call when the first exceeds a
        delay $h$; with $q = \Pr[T > h]$, expected load per request rises
        to $1 + q$ while the tail (e.g.\ $p99$) falls sharply. Constraints:
        hedge only idempotent operations (or carry idempotency keys),
        disregard the loser, and never hedge into an open circuit ---
        hedging a sick dependency is amplification, not resilience.
  \item \textbf{Client-side versus server-side load balancing?}
        Server-side: a dedicated LB owns health checking and instance
        selection; clients stay simple; cost is an extra hop and a
        component to make redundant. Client-side: every client watches a
        registry and picks instances; fewer hops, richer per-client policy;
        cost is embedding correct LB logic into every client in every
        language. Meshes move client-side logic into the sidecar to get
        both.
  \item \textbf{Explain the outbox pattern and its delivery guarantee.}
        Domain write and event-record insert commit in one local
        transaction (atomicity); a relay polls the outbox, publishes, and
        marks published only after the broker accepts. A crash between
        publish and mark replays the event, so the guarantee is
        \textbf{at-least-once} --- which is precisely why consumers must be
        idempotent. Together they yield effectively-once processing without
        distributed transactions (Listing~\ref{lst:ch15-outbox}).
  \item \textbf{What is backpressure, and name three concrete mechanisms.}
        The ability of an overloaded component to push load upstream
        instead of collapsing: explicit rejection (429/503 with
        \texttt{Retry-After}), bounded queues that fail fast when full, and
        credit-based flow control (HTTP/2 stream windows, gRPC). Unbounded
        queues are not backpressure; they are delayed failure.
  \item \textbf{Why is a slow dependency more dangerous than a dead one?}
        A dead dependency fails fast: errors are immediate and cheap, the
        breaker opens quickly, callers degrade. A slow dependency fails
        \emph{expensively}: every call occupies a worker, a connection, and
        memory for the full timeout, so the caller's pools exhaust ---
        stages 1--3 of Figure~\ref{fig:ch15-cascade} --- and the failure
        propagates to healthy functionality sharing those pools.
\end{enumerate}

\subsection*{Senior (Q21--Q27)}
\begin{enumerate}[label=\textbf{Q\arabic*.}, leftmargin=1.6cm, start=21]
  \item \textbf{Walk through a cascading failure from a latency blip to a
        fleet outage, and name the intervention at each stage.}
        Latency blip $\to$ blocked workers (intervene: per-hop timeouts)
        $\to$ pool exhaustion (bulkheads, bounded concurrency) $\to$ queue
        buildup (bounded queues, load shedding) $\to$ timeouts plus retries
        (budgets, jitter, breakers) $\to$ death spiral with flapping
        instances (gradual re-add, traffic ramps, staged restarts). Each
        stage's signal and countermeasure:
        Table~\ref{tab:ch15-failure-modes}.
  \item \textbf{Saga versus 2PC --- when and why?}
        2PC gives isolation at the price of blocking participants, coupled
        availability, and a coordinator whose crash leaves everyone
        in-doubt; over stateless HTTP it is unimplementable correctly.
        Sagas give atomicity of \emph{outcome} via local transactions plus
        compensating actions, accepting visible intermediate states
        (no isolation). Choose sagas for cross-service business flows in
        always-available systems --- which is to say, almost always; choose
        2PC essentially only inside a single tightly administered database
        tier~\cite{kleppmann2017ddia}.
  \item \textbf{Choreographed versus orchestrated saga --- trade-offs?}
        Choreography keeps each participant autonomous and removes the
        orchestrator as a bottleneck, but the workflow --- and worse, the
        compensation graph --- exists nowhere in code; debugging
        ``where is order ORD-7?'' requires reconstructing it from event
        logs. Orchestration centralizes sequence, state, and compensation,
        making status queries and failure handling explicit, at the cost of
        operating one more stateful component. Beyond two or three steps,
        orchestration is usually the honest choice.
  \item \textbf{What does a service mesh actually solve --- and what does
        it not solve?}
        Solves, uniformly and language-independently: mTLS workload
        identity, transport encryption, discovery, load balancing, per-hop
        retries/timeouts/circuit breaking, and golden-signal telemetry ---
        moved from application libraries into the data plane. Does not
        solve: idempotency of your endpoints, saga compensation logic,
        schema evolution, end-to-end deadline budgets across business
        operations, or knowing which of your failures are retryable. The
        mesh owns transport hygiene; the application still owns semantics.
  \item \textbf{How do you present eventual consistency honestly in an API
        contract?}
        Return \texttt{202 Accepted} with \texttt{Location} to a status
        resource; model the workflow's state machine as explicit,
        documented resource states (\texttt{PLACED}, \texttt{PAID},
        \texttt{SHIPPED}, \texttt{COMPENSATED}) with machine-readable
        failure reasons; give clients poll or webhook completion signals;
        never report a terminal state you have not reached. Consistency
        boundaries hidden from clients become their incidents.
  \item \textbf{Your breakers keep flapping open-closed-open after
        recovery. Diagnose.}
        Classic causes: $s = 1$ (one lucky probe closes onto a still-fragile
        dependency, re-fail, reopen), $d$ shorter than the dependency's
        real recovery time, and no traffic ramp --- the closed breaker
        admits a thundering herd of pent-up traffic. Fixes: require several
        consecutive probe successes, lengthen $d$ past measured recovery,
        and shed or ramp traffic on close. Alert on the flap rate itself.
  \item \textbf{How do bulkheads interact with async event loops?}
        An event loop makes threads cheap but does not make \emph{the
        dependency's} capacity infinite: without semaphores per dependency
        you can still issue ten thousand concurrent calls into a service
        that can serve two hundred, buying you its collapse. The async
        bulkhead is a semaphore (plus per-dependency connection limits);
        Listing~\ref{lst:ch15-demo} uses both. Concurrency is cheaper, not
        free.
\end{enumerate}

\subsection*{Staff (Q28--Q32)}
\begin{enumerate}[label=\textbf{Q\arabic*.}, leftmargin=1.6cm, start=28]
  \item \textbf{Design the resilience posture for a new user-facing
        endpoint with four downstream dependencies, two of which are
        non-critical.}
        Critical dependencies: deadline budgets derived from the endpoint
        SLO, one breaker each with thresholds justified from their SLOs and
        volumes, budgeted and jittered retries on idempotent calls only,
        per-dependency bulkheads sized below their healthy capacity.
        Non-critical: same machinery, plus designed degradation --- the
        endpoint renders without them (cached or partial responses), their
        budget share is small, and their breakers' open state is a silent
        fallback, not a 503. Then: a chaos drill proving the posture before
        launch, metrics hooks wired from day one, and an explicit answer
        for every dependency to ``what does the user get when this is
        down?''
  \item \textbf{When would you deliberately choose \emph{not} to retry?}
        When the operation is non-idempotent and carries no idempotency
        key; when the caller's remaining deadline cannot absorb the backoff
        (retrying produces work nobody awaits); when the error is
        deterministic (4xx validation, auth); when the budget is exhausted
        and the marginal retry's expected value is below its contribution
        to a developing retry storm; and in the middle of an incident,
        fleet-wide, where the kindest thing a client can do is fail fast
        and quietly.
  \item \textbf{How would you set organization-wide defaults for timeouts,
        retries, and breakers without creating false confidence?}
        Ship a thin client library encoding the chapter's machinery ---
        budgets, jitter, per-dependency breakers, metrics hooks --- with
        defaults derived from measured platform latency distributions, not
        folklore. Require per-service overrides to carry a justification
        and an expiration. Instrument everything and publish the
        flapping/open-rate dashboards. And run shared chaos drills,
        because a default nobody has seen fire is decoration
        (Section~\ref{sec:ch15-lab}). Policy encodes the floor; drills
        prove the behavior.
  \item \textbf{A 200 ms latency blip in one leaf service took down three
        tiers. Reconstruct the mechanism and the minimal fix set.}
        Theorem~\ref{thm:ch15-retry-amplification} with synchronized
        backoffs: the blip converted slow successes into timeouts; retries
        multiplied load $\approx a^{n}$; the leaf's queues deepened,
        raising latency further (the feedback arrow of
        Figure~\ref{fig:ch15-cascade}); shared pools carried the failure
        sideways to healthy endpoints. Minimal fixes, in order of leverage:
        per-hop deadline budgets; full jitter; process-wide retry budgets;
        per-dependency breakers and bulkheads; gradual recovery (probes,
        traffic ramps). Any one helps; the composition is the fix, because
        the outage was a property of the composition.
  \item \textbf{Design the migration from in-code resilience libraries to
        a service mesh for a forty-service platform --- what moves, what
        stays, and how do you avoid a regression mid-migration?}
        Moves to the mesh: mTLS and identity, discovery, load balancing,
        per-hop timeout/retry/breaker enforcement, telemetry. Stays in the
        application: idempotency, saga orchestration and compensation,
        deadline budgets across business operations, degradation/fallback
        semantics, retryability classification. Migration discipline:
        shadow first (mesh enforcing telemetry only), then enforce with
        in-code policy as the safety net --- and make the two policies
        \emph{identical} during overlap, because mesh retries stacked on
        library retries are amplification you chose on purpose. Remove
        in-code transport policy only after drills against the mesh
        configuration reproduce Table~\ref{tab:ch15-drill}'s behavior.
        Treat it as a protocol migration
        (Chapter~\ref{ch:versioning}'s discipline), not an install.
\end{enumerate}

%% ----------------------------------------------------------------------------
\section{External Bibliography \& Further Reading}
\label{sec:ch15-bibliography}

\begin{itemize}
  \item Sam Newman, \emph{Building Microservices}, 2nd ed., O'Reilly, 2021
        --- the reference for communication styles, orchestration versus
        choreography, and sagas in organizational
        context~\cite{newman2021microservices}.
  \item Michael T. Nygard, \emph{Release It!}, 2nd ed., Pragmatic
        Bookshelf, 2018 --- the origin text for circuit breakers,
        bulkheads, timeouts, and the stability
        patterns/anti-patterns vocabulary this chapter
        formalizes~\cite{nygard2018releaseit}.
  \item Martin Kleppmann, \emph{Designing Data-Intensive Applications},
        O'Reilly, 2017 --- chapters on distributed systems faults,
        consistency, and exactly-why-2PC-is-hard; the theoretical
        foundation for Section~\ref{subsec:ch15-transactions}
        ~\cite{kleppmann2017ddia}.
  \item Uwe Friedrichsen, \emph{Patterns of Resilience} article series
        (uwe-friedrichsen.de) --- a deep, pragmatic taxonomy of resilience
        patterns and their combinations; the best long-form companion to
        Sections~\ref{subsec:ch15-timeouts}--\ref{subsec:ch15-cascade}.
  \item Betsy Beyer et al., \emph{Site Reliability Engineering}, O'Reilly,
        2016 --- the chapters ``Addressing Cascading Failures'' and
        ``Managing Load'' (overload, load shedding, criticality) at Google
        scale.
  \item Gregor Hohpe and Bobby Woolf, \emph{Enterprise Integration
        Patterns}, Addison-Wesley, 2003 --- the canonical vocabulary for
        messaging, idempotent receivers, and the patterns underlying
        Section~\ref{subsec:ch15-transactions}.
  \item Envoy Proxy documentation (envoyproxy.io) --- circuit breaking,
        outlier detection, retries, and the xDS APIs the mesh figure
        (Figure~\ref{fig:ch15-mesh}) abstracts.
  \item Istio documentation (istio.io) --- control-plane architecture,
        traffic management, and mTLS identity model.
  \item The gRPC documentation on deadlines and
        cancellation~\cite{grpcdocs}; the FastAPI~\cite{fastapidocs} and
        httpx~\cite{httpxdocs} documentation for the toolkit used
        throughout the listings; and the twelve-factor
        methodology~\cite{twelvefactor} for the operational posture this
        chapter's services assume.
\end{itemize}

%% ----------------------------------------------------------------------------
\section{Capstone Projects}
\label{sec:ch15-capstones}

\subsection*{Capstone 15.1 --- Breaker Under a Microscope \hfill [junior]}
\textbf{Objective:} Reimplement Listing~\ref{lst:ch15-breaker} from
Definition~\ref{def:ch15-breaker} alone, then attack it.

\textbf{Inputs/Datasets:} Listings 15.1--15.2; no datasets.

\textbf{Steps:} (1)~Reimplement the breaker from the definition without
consulting the listing. (2)~Run Listing~\ref{lst:ch15-breaker-tests}
unmodified against your implementation. (3)~Seed three deliberate bugs ---
drop the minimum-call floor, forget window eviction, allow unlimited
probes --- and record which tests catch each. (4)~Add one new test:
\texttt{half\_open\_success\_threshold=3} with a probe sequence of
success--success--failure, asserting the breaker reopens and requires the
full open duration again.

\textbf{Acceptance Criteria:}
\begin{itemize}
  \item[$\square$] All Listing~\ref{lst:ch15-breaker-tests} tests pass
        against your reimplementation.
  \item[$\square$] Each seeded bug is caught by at least one existing test;
        you can name the mapping.
  \item[$\square$] The new probe-threshold test passes and fails against
        the seeded unlimited-probe bug.
\end{itemize}

\textbf{Stretch Goals:} Add the slow-call failure dimension (Exercise 4 of
Section~\ref{sec:ch15-exercises}) with a manual-clock proof.

\subsection*{Capstone 15.2 --- Latency Chaos and the Price of Hedging \hfill [mid]}
\textbf{Objective:} Extend the drill to latency faults and measure hedged
requests honestly.

\textbf{Inputs/Datasets:} Listings 15.5--15.6; the \texttt{ChaosConfig}
latency knob.

\textbf{Steps:} (1)~Drive inventory with \texttt{latency\_seconds} drawn
from a seeded bimodal distribution ($95\%$ fast, $5\%$ slow). (2)~Measure
catalog $p99$ with no hedging. (3)~Add hedging at the measured fast-mode
$p95$ and re-measure $p99$ and the extra-call rate. (4)~Verify the extra
load matches $1 + q$ within sampling error across seeded runs.

\textbf{Acceptance Criteria:}
\begin{itemize}
  \item[$\square$] Hedged $p99$ is materially below unhedged $p99$ on the
        same seed.
  \item[$\square$] Measured hedge rate equals $1 + q$ within the tolerance
        you derived.
  \item[$\square$] Hedging is disabled while the breaker is open; a test
        proves it.
  \item[$\square$] Everything is seeded and offline; no wall-clock
        assertions.
\end{itemize}

\textbf{Stretch Goals:} Cancel the loser request and measure the downstream
savings; add a budget so hedges spend retry tokens too.

\subsection*{Capstone 15.3 --- Production Outbox: Retry, DLQ, Ordering \hfill [mid]}
\textbf{Objective:} Harden Listing~\ref{lst:ch15-outbox} into the relay you
would actually run.

\textbf{Inputs/Datasets:} Listings 15.7--15.8; a seeded event generator you
write (Northwind Robotics orders, 10k events).

\textbf{Steps:} (1)~Add relay retry with the
Listing~\ref{lst:ch15-retry-budget} budget for publish failures.
(2)~Add a dead-letter column/table after $k$ failed attempts, with a
requeue operation. (3)~Enforce per-aggregate ordering: no event for
aggregate $A$ is published while an earlier event for $A$ is unpublished
after a failure --- implement via per-aggregate sequencing in the outbox.
(4)~Tests: poison event goes to DLQ after exactly $k$ attempts; aggregates
without the poison event are unaffected; requeue restores flow exactly
once.

\textbf{Acceptance Criteria:}
\begin{itemize}
  \item[$\square$] Publish-failure retries respect the budget cap (proved
        with a counting publisher).
  \item[$\square$] DLQ after exactly $k$ attempts; requeue is idempotent.
  \item[$\square$] Per-aggregate ordering holds under crash/replay
        injection; cross-aggregate throughput is unaffected.
  \item[$\square$] The Listing~\ref{lst:ch15-outbox-tests} suite still
        passes unmodified.
\end{itemize}

\textbf{Stretch Goals:} Add a lag metric (unpublished count, oldest
unpublished age) to the relay's event hook and alert thresholds to the
lab's dashboard list.

\subsection*{Capstone 15.4 --- Measure Theorem~\ref{thm:ch15-retry-amplification} \hfill [senior]}
\textbf{Objective:} Build a four-hop amplification testbed and verify the
theorem empirically, then deploy the countermeasures and verify them too.

\textbf{Inputs/Datasets:} Listings 15.1, 15.3, 15.5; a fourth FastAPI tier
you add downstream of warehouse.

\textbf{Steps:} (1)~Wire edge $\to$ catalog $\to$ inventory $\to$
warehouse, each hop with a counting stub for its downstream and
configurable retry count. (2)~With the leaf hard-down and three attempts
per hop, measure calls arriving at each hop; compare to $3, 9, 27$ and the
total to $39$. (3)~Add per-process one-token retry budgets at the two
middle hops; re-measure; derive the closed form for the budgeted case and
check it. (4)~Add full jitter and confirm (statistically, over seeds) that
retry burst correlation drops.

\textbf{Acceptance Criteria:}
\begin{itemize}
  \item[$\square$] Measured per-hop counts match $a^k$ exactly in the
        unbudgeted run.
  \item[$\square$] Your closed-form prediction for the budgeted run matches
        measurement.
  \item[$\square$] A written paragraph explains any residual gap between
        theory and measurement.
  \item[$\square$] Fully offline, seeded, no wall-clock assertions.
\end{itemize}

\textbf{Stretch Goals:} Replace counting stubs with the real ASGI chain and
show the theorem survives real HTTP machinery.

\subsection*{Capstone 15.5 --- Resilience Architecture for Fleet Telemetry \hfill [staff]}
\textbf{Objective:} Produce and defend a complete resilience design for
Northwind Robotics' fleet-telemetry platform: ingestion API $\to$
normalizer $\to$ enrichment $\to$ storage, plus a user-facing dashboard
read path.

\textbf{Inputs/Datasets:} This chapter; the Northwind Robotics endpoint set
from prior chapters; SLO sheet you write (ingestion $99.95\%$,
dashboard $99.9\%$ at $p99 < 800\,\mathrm{ms}$).

\textbf{Steps:} (1)~Draw the call graph and classify every edge: sync/async,
coupling dimensions, retryability, criticality. (2)~Assign deadline
budgets, breaker parameters, bulkhead sizes, retry budgets, and
degradation behavior per edge, each justified from the SLOs and a stated
traffic model. (3)~Write the failure-mode analysis table (extend
Table~\ref{tab:ch15-failure-modes}) for your graph's top five scenarios,
including one poisoning attack and one slow-dependency case. (4)~Define
the chaos-drill program: which faults, which cadence, which abort
criteria. (5)~Record the mesh boundary: which policies move to the mesh,
which stay in the application, and why.

\textbf{Acceptance Criteria:}
\begin{itemize}
  \item[$\square$] Every edge carries a justified, SLO-derived parameter
        set; no ``default because default.''
  \item[$\square$] Deadline budgets satisfy
        Definition~\ref{def:ch15-deadline} on every path.
  \item[$\square$] The failure-mode table has a signal and a countermeasure
        per scenario; at least one scenario requires degradation, not just
        failure.
  \item[$\square$] The drill program includes abort criteria and a
        rollback story.
  \item[$\square$] The mesh boundary list contains at least three
        application-owned concerns with correct justifications.
\end{itemize}

\textbf{Stretch Goals:} Cost the design: extra requests from retries and
hedges, engineering hours, and drill operational cost, against the
incident cost model from Section~\ref{sec:ch15-lab}'s war story.

%% ----------------------------------------------------------------------------
\section{Python Dependencies Appendix}
\label{sec:ch15-deps}

All code runs offline after a one-time install. From PowerShell:

\begin{minted}{text}
python -m venv .venv
.\.venv\Scripts\Activate.ps1
pip install -r requirements.txt
\end{minted}

\noindent\textbf{requirements.txt:}
\begin{minted}{text}
fastapi>=0.110
httpx>=0.27
pydantic>=2.7
pytest>=8
tenacity>=8.3
\end{minted}

\subsection{fastapi ($\geq$ 0.110)}
\textbf{Description:} The ASGI framework hosting the catalog, inventory,
and warehouse services of Listing~\ref{lst:ch15-demo}. \textbf{Why
included:} real routing, dependency injection, and response handling over
real ASGI request cycles --- the chaos drill exercises genuine framework
machinery, not a mock of it~\cite{fastapidocs}. \textbf{Alternatives:}
Starlette directly (the demo would work nearly unchanged); Flask (WSGI ---
loses the async concurrency the drill depends on); a live multi-process
deployment under uvicorn (faithful, but no longer offline-pure).
\textbf{Installation:} \texttt{pip install "fastapi>=0.110"} in the
activated virtual environment.

\subsection{httpx ($\geq$ 0.27)}
\textbf{Description:} The sync/async HTTP client used for every
inter-service call in this chapter~\cite{httpxdocs}. \textbf{Why included:}
\texttt{ASGITransport} wires clients to in-process FastAPI apps with the
real request/response and timeout code paths --- no sockets, no network;
\texttt{Limits} provides the connection-pool bulkhead; the
\texttt{TransportError} hierarchy is what
Listing~\ref{lst:ch15-retry-budget} classifies. \textbf{Alternatives:}
\texttt{requests} (sync only, no ASGI transport); \texttt{aiohttp} (fine,
but its client/timeout model differs from this book's other chapters);
gRPC with its own deadline propagation for the RPC variant
(Chapter~\ref{ch:grpc}). \textbf{Installation:}
\texttt{pip install "httpx>=0.27"}.

\subsection{pydantic ($\geq$ 2.7)}
\textbf{Description:} Data validation and settings library; the model layer
under FastAPI. \textbf{Why included:} it is FastAPI's contract surface ---
request/response schemas, validation errors --- and the natural home for
the typed configuration objects (breaker and retry policies) you will
extract when these listings grow configuration files. \textbf{Alternatives:}
\texttt{dataclasses} with hand-rolled validation (what the listings use for
their internal value types, to keep the dependency surface visible);
\texttt{attrs}; \texttt{msgspec} (fast, narrower ecosystem).
\textbf{Installation:} \texttt{pip install "pydantic>=2.7"}.

\subsection{pytest ($\geq$ 8)}
\textbf{Description:} The test framework running all 34 chapter tests.
\textbf{Why included:} fixtures, \texttt{raises}, and \texttt{approx} keep
the deterministic proofs (manual clocks, counting stubs, event-order
assertions) short enough to read as specifications. \textbf{Alternatives:}
\texttt{unittest} (stdlib; workable, more boilerplate); \texttt{nose2}
(unmaintained momentum). \textbf{Installation:}
\texttt{pip install "pytest>=8"}.

\subsection{tenacity ($\geq$ 8.3)}
\textbf{Description:} The production-grade Python retry library:
declarative stop/wait/retry conditions, jittered exponential backoff,
async support. \textbf{Why included:} the chapter's
Listing~\ref{lst:ch15-retry-budget} deliberately teaches the machinery by
building it; in real services you should reach for tenacity's
battle-tested implementation and add this chapter's two contributions ---
the process-wide token-bucket budget and the deadline guard --- on top.
\textbf{Alternatives:} \texttt{pybreaker} for the circuit-breaker half
(mature, sync-first; compare its state machine against
Definition~\ref{def:ch15-breaker} as an exercise); \texttt{resilience4j}
concepts if your estate is JVM (the patterns are identical; the Java
implementation is the genre's reference); mesh-level retry policy
(Envoy/Istio) for the transport tier --- with
Section~\ref{sec:ch15-pitfalls} P15.7's caveat. \textbf{Installation:}
\texttt{pip install "tenacity>=8.3"}.

%% ----------------------------------------------------------------------------
\section{Chapter Summary \& Key Takeaways}
\label{sec:ch15-summary}

\begin{itemize}
  \item Distribution moves coupling from the compiler to the network.
        Classify every interaction by sync/async and
        orchestration/choreography, and by its temporal, spatial, and
        logical coupling --- then make each coupling a decision, not a
        default (Section~\ref{subsec:ch15-taxonomy}).
  \item The fallacies of distributed computing are a checklist, not a
        trivia list; each ignored fallacy has a named failure mode, an
        early signal, and a metric (Section~\ref{subsec:ch15-coupling}).
  \item Timeouts are the foundation: every call gets one, and budgets
        shrink strictly toward the leaves
        (Definition~\ref{def:ch15-deadline}). A timeout longer than the
        caller's manufactures wasted work.
  \item Retries amplify multiplicatively: $S(n,a) = a(a^n-1)/(a-1)$
        downstream requests per client request --- $39\times$ for three
        attempts across three hops
        (Theorem~\ref{thm:ch15-retry-amplification}). Token-bucket budgets
        cap it (Proposition~\ref{prop:ch15-budget-cap}); jitter decorrelates
        it; deadlines and \texttt{Retry-After} discipline it.
  \item The circuit breaker is a state machine with parameters
        $(\theta, W, N, d, p, s)$, one per dependency: minimum-call floors
        against noise, rolling windows against staleness, bounded probes
        against stampedes (Definition~\ref{def:ch15-breaker},
        Listing~\ref{lst:ch15-breaker}).
  \item Bulkheads bound failure domains; hedged requests buy tail latency
        at an explicit $1 + q$ load tax; backpressure converts overload
        into cheap rejection instead of expensive collapse.
  \item Cascading failure is a six-stage spiral with a feedback loop; every
        stage has an observable precursor and a countermeasure
        (Figure~\ref{fig:ch15-cascade},
        Table~\ref{tab:ch15-failure-modes}). Recovery must be engineered:
        probes, ramps, and staged restarts, because a fallen service
        cannot stand back up into the load that felled it.
  \item Discovery and load balancing decide \emph{who talks to whom};
        the mesh moves transport hygiene (mTLS, retries, timeouts,
        telemetry) into the data plane while the control plane configures
        it --- and neither touches application semantics, which remain
        yours (Section~\ref{subsec:ch15-mesh}).
  \item Distributed transactions: sagas with compensating actions instead
        of 2PC; the transactional outbox for atomic write-plus-publish;
        idempotent consumers to absorb at-least-once replays; and eventual
        consistency presented honestly in the API contract.
  \item Patterns you have not drilled are decoration. The chaos harness
        exists so you can watch the state machine fire before the incident
        does (Section~\ref{sec:ch15-lab}).
\end{itemize}

%% ----------------------------------------------------------------------------
\section{Quick-Reference Card}
\label{sec:ch15-quickref}

\subsection*{Pattern selection}
\begin{table}[h]
  \centering
  \small
  \begin{tabularx}{\textwidth}{l X}
    \toprule
    \textbf{Symptom / risk} & \textbf{Reach for} \\
    \midrule
    Calls hang when dependency slows & per-hop timeouts + deadline budgets \\
    Retry volume exploding during incidents & token-bucket retry budget + full jitter \\
    One dependency poisoning the process & circuit breaker (one per dependency) + bulkhead \\
    Tail latency dominates user experience & hedged requests (idempotent calls only) \\
    Overload converting to collapse & backpressure: bounded queues, 429/503 + \texttt{Retry-After} \\
    Dual-write inconsistency (DB + events) & transactional outbox + relay \\
    Replayed events double-apply effects & idempotent consumer (dedup in same tx) \\
    Multi-step flow fails halfway & saga with compensating actions (orchestrated) \\
    Polyglot transport policy drift & service mesh data plane (keep semantics in-app) \\
    \bottomrule
  \end{tabularx}
  \caption{Symptom-to-pattern selection.}
  \label{tab:ch15-pattern-select}
\end{table}

\subsection*{Breaker parameter cheat-sheet}
\begin{table}[h]
  \centering
  \small
  \begin{tabularx}{\textwidth}{l l X}
    \toprule
    \textbf{Parameter} & \textbf{Start at} & \textbf{Tune by} \\
    \midrule
    $\theta$ failure rate & 0.5 & lower for flaky-but-critical, never below SLO noise \\
    $W$ window & 10--60 s & dependency's failure-burst duration \\
    $N$ minimum calls & 10--20 & high-volume deps: higher; low-volume: use duration windows \\
    $d$ open duration & 5--30 s & measured recovery time of the dependency \\
    $p$ probe concurrency & 1 & raise only with a traffic-ramp story \\
    $s$ probe successes & 2--3 & raise if you observe flap (open/close oscillation) \\
    \bottomrule
  \end{tabularx}
  \caption{Starting points for Definition~\ref{def:ch15-breaker}'s
  parameters. Every number is a hypothesis until a chaos drill confirms it.}
  \label{tab:ch15-breaker-cheatsheet}
\end{table}

\subsection*{Anti-amplification checklist}
\begin{itemize}
  \item[$\square$] Every outbound call has a timeout; budgets shrink toward
        the leaves on every path.
  \item[$\square$] Retries only on idempotent operations or with
        idempotency keys (Chapter~\ref{ch:idempotency}).
  \item[$\square$] Every retry spends a token from a process-wide budget;
        first attempts are free.
  \item[$\square$] Backoff is exponential with full jitter; no fixed
        sleeps anywhere in the fleet.
  \item[$\square$] Call-level deadlines bound total retry time; overshoot
        skips the retry.
  \item[$\square$] \texttt{Retry-After} from 429/503 responses overrides
        local schedules.
  \item[$\square$] One breaker per dependency; bulkheads sized below each
        dependency's healthy capacity.
  \item[$\square$] Breaker and budget events flow to metrics; opening and
        flapping are alertable.
  \item[$\square$] The chaos drill runs on a schedule, with abort criteria
        --- and someone watches the state transitions.
\end{itemize}

\section{Standardized Evidence Addendum}
\subsection{Benchmark Snapshot Template}
\begin{table}[h]
  \centering
  \small
  \begin{tabularx}{\textwidth}{l c c c c}
    \toprule
    \textbf{Config} & \textbf{p50 latency} & \textbf{p99 latency} & \textbf{Error rate} & \textbf{Throughput} \\
    \midrule
    Baseline & -- & -- & -- & 1.00x \\
    Candidate A & -- & -- & -- & -- \\
    Candidate B & -- & -- & -- & -- \\
    \bottomrule
  \end{tabularx}
  \caption{Minimum benchmark table to keep per chapter.}
\end{table}

\subsection{Request Trace Template}
\begin{itemize}
  \item \textbf{Request:} one representative API call for this chapter's topic.
  \item \textbf{Before:} behavior (headers, payload, latency, failure mode) with the naive approach.
  \item \textbf{After:} behavior with the technique in this chapter.
  \item \textbf{Result:} what changed in correctness, resilience, latency, and operability.
\end{itemize}

\subsection{Cost and Latency Budget Snapshot}
\begin{table}[h]
  \centering
  \small
  \begin{tabularx}{\textwidth}{l c c}
    \toprule
    \textbf{Stage} & \textbf{Latency budget} & \textbf{Cost driver} \\
    \midrule
    TLS / connection setup & -- & handshake round trips, keep-alive hit rate \\
    Request validation & -- & schema complexity, CPU per request \\
    Business logic and I/O & -- & downstream fan-out, query cost \\
    Serialization and egress & -- & payload size, compression ratio \\
    \bottomrule
  \end{tabularx}
  \caption{Operational budget template for this chapter's technique.}
\end{table}

\section{Configuration Preset Template}
\begin{minted}{text}
# api_policy.yaml
policy_version: "v1.0.0"
component: "<chapter_topic>"
timeout_ms: 0
retry_budget: 0
fallback_strategy: "fail_fast"
rollback_trigger: "error_budget_burn_gt_threshold"
\end{minted}

\section{CI Quality Gate Specification}
\begin{itemize}
  \item contract tests green against the pinned OpenAPI/AsyncAPI/Protobuf spec,
  \item no breaking schema changes without a version bump,
  \item error responses conform to RFC 9457 problem-details shape,
  \item benchmark deltas within approved regression bounds (latency and throughput),
  \item policy version and rollback plan present in release metadata.
\end{itemize}

\section{Incident Runbook Template}
\begin{enumerate}
  \item Detect regression from SLO burn-rate alerts or consumer reports.
  \item Scope impacted endpoints, versions, and consumer segments.
  \item Contain impact (rollback, feature flag, rate-limit tightening, or traffic shift).
  \item Diagnose with traces, structured logs, and contract diffs.
  \item Recover with validated fix and verified rollout procedure.
  \item Prevent recurrence via new regression tests and CI gates.
\end{enumerate}

\section{Cross-Chapter Decision Card}
\begin{table}[h]
  \centering
  \small
  \begin{tabularx}{\textwidth}{l X X}
    \toprule
    \textbf{Dimension} & \textbf{Choose this approach when} & \textbf{Prefer an alternative when} \\
    \midrule
    Use case & -- & -- \\
    Latency budget & -- & -- \\
    Operational cost & -- & -- \\
    Team maturity & -- & -- \\
    \bottomrule
  \end{tabularx}
  \caption{Decision card to complete per chapter.}
\end{table}

\section{ROI Model and Build-vs-Buy}
\begin{itemize}
  \item \textbf{Build when:} the capability is differentiating, compliance forbids vendors, or scale makes vendor pricing irrational.
  \item \textbf{Buy when:} the capability is commodity (auth, gateways, observability), time-to-market dominates, or staffing is thin.
  \item \textbf{Hidden costs to model:} on-call load, upgrade cadence, expertise ramp, data egress fees.
\end{itemize}

\section{Disaster Recovery Notes (RPO/RTO)}
\begin{itemize}
  \item Define recovery point objective (RPO): maximum acceptable data loss window.
  \item Define recovery time objective (RTO): maximum acceptable restoration time.
  \item Test restores regularly; an untested backup is not a backup.
\end{itemize}

